% =============================================================
%  preamble.tex — Reusable preamble for Neural Networks course
%  Usage: % =============================================================
%  preamble.tex — Reusable preamble for Neural Networks course
%  Usage: % =============================================================
%  preamble.tex — Reusable preamble for Neural Networks course
%  Usage: % =============================================================
%  preamble.tex — Reusable preamble for Neural Networks course
%  Usage: \input{preamble} at the top of your document,
%         after \documentclass but before \begin{document}.
% =============================================================

% --- Encoding & language ---
\usepackage[T1]{fontenc}
\usepackage[utf8]{inputenc}
\usepackage[english]{babel}         % swap for [french] if needed

% --- Page geometry ---
\usepackage{geometry}
\geometry{top=2.5cm, bottom=2.5cm, left=2.5cm, right=2.5cm}

% --- Math ---
\usepackage{amsmath, amssymb}

% --- Tables ---
\usepackage{booktabs}               % \toprule, \midrule, \bottomrule
\usepackage{array}                  % column type extensions
\usepackage{tabularx}               % auto-width columns
\usepackage{multirow}               % multi-row cells

% --- Colors ---
\usepackage{xcolor}

\definecolor{mainblue}{RGB}{30, 80, 160}
\definecolor{lightblue}{RGB}{220, 235, 255}
\definecolor{accentgray}{RGB}{245, 246, 248}
\definecolor{darkgray}{RGB}{80, 80, 80}
\definecolor{solutiongreen}{RGB}{20, 120, 60}
\definecolor{lightgreen}{RGB}{220, 245, 230}
\definecolor{warnorange}{RGB}{210, 120, 20}
\definecolor{lightorange}{RGB}{255, 243, 220}

% --- Lists ---
\usepackage{enumitem}

% --- TikZ & pgfplots ---
\usepackage{tikz}
\usepackage{pgfplots}
\pgfplotsset{compat=1.18}
\usetikzlibrary{positioning, arrows.meta, calc, patterns}

% --- Boxes: tcolorbox ---
\usepackage{tcolorbox}
\tcbuselibrary{skins, breakable}

% Generic exercise box (blue border, gray background)
\tcbset{
  exercisebox/.style={
    enhanced,
    colback=accentgray,
    colframe=mainblue,
    boxrule=0.8pt,
    arc=4pt,
    left=10pt, right=10pt, top=8pt, bottom=8pt,
    fonttitle=\bfseries\color{mainblue},
  }
}

% Solution box (green, used in answer sheets)
\tcbset{
  solutionbox/.style={
    enhanced,
    colback=lightgreen,
    colframe=solutiongreen,
    boxrule=0.8pt,
    arc=4pt,
    left=10pt, right=10pt, top=8pt, bottom=8pt,
    fonttitle=\bfseries\color{solutiongreen},
  }
}

% Answer highlight box (orange, for final numeric answers)
\tcbset{
  answerbox/.style={
    enhanced,
    colback=lightorange,
    colframe=warnorange,
    boxrule=0.7pt,
    arc=4pt,
    left=8pt, right=8pt, top=5pt, bottom=5pt,
    fonttitle=\bfseries\color{warnorange},
  }
}

% Step box (light blue border, gray background, for intermediate steps)
\tcbset{
  stepbox/.style={
    enhanced,
    colback=accentgray,
    colframe=mainblue!40,
    boxrule=0.5pt,
    arc=3pt,
    left=8pt, right=8pt, top=6pt, bottom=6pt,
  }
}

% --- Inline note / warning box (mdframed) ---
\usepackage{mdframed}
\newmdenv[
  linecolor=warnorange,
  backgroundcolor=orange!5,
  roundcorner=4pt,
  innerleftmargin=8pt,
  innerrightmargin=8pt,
  innertopmargin=6pt,
  innerbottommargin=6pt,
  linewidth=0.7pt
]{notebox}

% --- Section formatting ---
\usepackage{titlesec}

% Section style: "Exercise N : Title" with a blue rule underneath
% Change the prefix via \renewcommand{\sectionprefix}{...} before \section
\newcommand{\sectionprefix}{Exercise}
\titleformat{\section}[block]
  {\large\bfseries\color{mainblue}}
  {\sectionprefix~\thesection\,:}{0.6em}{}
  [\vspace{-0.3em}\textcolor{mainblue}{\rule{\linewidth}{0.8pt}}]
\titlespacing{\section}{0pt}{1.6em}{0.8em}

% --- Header / Footer ---
\usepackage{fancyhdr}
\pagestyle{fancy}
\fancyhf{}
\renewcommand{\headrulewidth}{0.4pt}

% Customise these three commands in your document after \input{preamble}:
%   \renewcommand{\doccoursename}{My Course}
%   \renewcommand{\docsessionname}{Session N — Topic}
\newcommand{\doccoursename}{Neural Networks for Engineers}
\newcommand{\docsessionname}{Session}

\fancyhead[L]{\small\textcolor{darkgray}{\doccoursename}}
\fancyhead[R]{\small\textcolor{darkgray}{\docsessionname}}
\fancyfoot[C]{\small\textcolor{darkgray}{\thepage}}

% --- Title block macro ---
% Usage: \maketitleblock{Title line}{Subtitle / metadata line}{background color}
% Examples:
%   \maketitleblock{Session 4 — Exercises}{Mixed level | Neural Networks}{mainblue}
%   \maketitleblock{Session 4 — Solutions}{Perceptrons and MLPs | Neural Networks}{solutiongreen}
\newcommand{\maketitleblock}[3]{%
  \begin{tcolorbox}[
    enhanced,
    colback=#3,
    colframe=#3,
    arc=6pt, boxrule=0pt,
    left=16pt, right=16pt, top=12pt, bottom=12pt
  ]
    {\large\bfseries\color{white} #1}\\[4pt]
    {\small\color{white!85!black} #2}
  \end{tcolorbox}
  \vspace{1em}
}

% --- Convenience macros ---

% Blank answer field (inline underline for fill-in tables)
\newcommand{\acell}{\rule{1.8cm}{0.4pt}}

% Step-by-step computation display: \step{label}{math}
% Example: \step{Hidden layer}{z^{(1)} = Wx + b = \ldots}
\newcommand{\step}[2]{%
  \noindent\textbf{#1:}\quad $#2$\par\vspace{0.3em}
}

% Neuron style for TikZ diagrams (reuse with \node[neuron])
% Colors are overridable per layer — see examples below.
%
% Input neuron  : fill=lightblue,  draw=mainblue
% Hidden neuron : fill=green!15,   draw=green!50!black
% Output neuron : fill=orange!15,  draw=orange!70!black
%
% Example architecture diagram:
%
%   \begin{tikzpicture}[
%     neuron/.style={circle, minimum size=1cm, thick, font=\small\bfseries},
%     arr/.style={->, >=Stealth, thick, gray!70}
%   ]
%     \node[neuron, fill=lightblue,  draw=mainblue]          (x1) at (0, 1)  {$x_1$};
%     \node[neuron, fill=green!15,   draw=green!50!black]    (h1) at (3, 1)  {$h_1$};
%     \node[neuron, fill=orange!15,  draw=orange!70!black]   (y)  at (6, 1)  {$\hat{y}$};
%     \draw[arr] (x1) -- (h1);
%     \draw[arr] (h1) -- (y);
%   \end{tikzpicture}
 at the top of your document,
%         after \documentclass but before \begin{document}.
% =============================================================

% --- Encoding & language ---
\usepackage[T1]{fontenc}
\usepackage[utf8]{inputenc}
\usepackage[english]{babel}         % swap for [french] if needed

% --- Page geometry ---
\usepackage{geometry}
\geometry{top=2.5cm, bottom=2.5cm, left=2.5cm, right=2.5cm}

% --- Math ---
\usepackage{amsmath, amssymb}

% --- Tables ---
\usepackage{booktabs}               % \toprule, \midrule, \bottomrule
\usepackage{array}                  % column type extensions
\usepackage{tabularx}               % auto-width columns
\usepackage{multirow}               % multi-row cells

% --- Colors ---
\usepackage{xcolor}

\definecolor{mainblue}{RGB}{30, 80, 160}
\definecolor{lightblue}{RGB}{220, 235, 255}
\definecolor{accentgray}{RGB}{245, 246, 248}
\definecolor{darkgray}{RGB}{80, 80, 80}
\definecolor{solutiongreen}{RGB}{20, 120, 60}
\definecolor{lightgreen}{RGB}{220, 245, 230}
\definecolor{warnorange}{RGB}{210, 120, 20}
\definecolor{lightorange}{RGB}{255, 243, 220}

% --- Lists ---
\usepackage{enumitem}

% --- TikZ & pgfplots ---
\usepackage{tikz}
\usepackage{pgfplots}
\pgfplotsset{compat=1.18}
\usetikzlibrary{positioning, arrows.meta, calc, patterns}

% --- Boxes: tcolorbox ---
\usepackage{tcolorbox}
\tcbuselibrary{skins, breakable}

% Generic exercise box (blue border, gray background)
\tcbset{
  exercisebox/.style={
    enhanced,
    colback=accentgray,
    colframe=mainblue,
    boxrule=0.8pt,
    arc=4pt,
    left=10pt, right=10pt, top=8pt, bottom=8pt,
    fonttitle=\bfseries\color{mainblue},
  }
}

% Solution box (green, used in answer sheets)
\tcbset{
  solutionbox/.style={
    enhanced,
    colback=lightgreen,
    colframe=solutiongreen,
    boxrule=0.8pt,
    arc=4pt,
    left=10pt, right=10pt, top=8pt, bottom=8pt,
    fonttitle=\bfseries\color{solutiongreen},
  }
}

% Answer highlight box (orange, for final numeric answers)
\tcbset{
  answerbox/.style={
    enhanced,
    colback=lightorange,
    colframe=warnorange,
    boxrule=0.7pt,
    arc=4pt,
    left=8pt, right=8pt, top=5pt, bottom=5pt,
    fonttitle=\bfseries\color{warnorange},
  }
}

% Step box (light blue border, gray background, for intermediate steps)
\tcbset{
  stepbox/.style={
    enhanced,
    colback=accentgray,
    colframe=mainblue!40,
    boxrule=0.5pt,
    arc=3pt,
    left=8pt, right=8pt, top=6pt, bottom=6pt,
  }
}

% --- Inline note / warning box (mdframed) ---
\usepackage{mdframed}
\newmdenv[
  linecolor=warnorange,
  backgroundcolor=orange!5,
  roundcorner=4pt,
  innerleftmargin=8pt,
  innerrightmargin=8pt,
  innertopmargin=6pt,
  innerbottommargin=6pt,
  linewidth=0.7pt
]{notebox}

% --- Section formatting ---
\usepackage{titlesec}

% Section style: "Exercise N : Title" with a blue rule underneath
% Change the prefix via \renewcommand{\sectionprefix}{...} before \section
\newcommand{\sectionprefix}{Exercise}
\titleformat{\section}[block]
  {\large\bfseries\color{mainblue}}
  {\sectionprefix~\thesection\,:}{0.6em}{}
  [\vspace{-0.3em}\textcolor{mainblue}{\rule{\linewidth}{0.8pt}}]
\titlespacing{\section}{0pt}{1.6em}{0.8em}

% --- Header / Footer ---
\usepackage{fancyhdr}
\pagestyle{fancy}
\fancyhf{}
\renewcommand{\headrulewidth}{0.4pt}

% Customise these three commands in your document after % =============================================================
%  preamble.tex — Reusable preamble for Neural Networks course
%  Usage: \input{preamble} at the top of your document,
%         after \documentclass but before \begin{document}.
% =============================================================

% --- Encoding & language ---
\usepackage[T1]{fontenc}
\usepackage[utf8]{inputenc}
\usepackage[english]{babel}         % swap for [french] if needed

% --- Page geometry ---
\usepackage{geometry}
\geometry{top=2.5cm, bottom=2.5cm, left=2.5cm, right=2.5cm}

% --- Math ---
\usepackage{amsmath, amssymb}

% --- Tables ---
\usepackage{booktabs}               % \toprule, \midrule, \bottomrule
\usepackage{array}                  % column type extensions
\usepackage{tabularx}               % auto-width columns
\usepackage{multirow}               % multi-row cells

% --- Colors ---
\usepackage{xcolor}

\definecolor{mainblue}{RGB}{30, 80, 160}
\definecolor{lightblue}{RGB}{220, 235, 255}
\definecolor{accentgray}{RGB}{245, 246, 248}
\definecolor{darkgray}{RGB}{80, 80, 80}
\definecolor{solutiongreen}{RGB}{20, 120, 60}
\definecolor{lightgreen}{RGB}{220, 245, 230}
\definecolor{warnorange}{RGB}{210, 120, 20}
\definecolor{lightorange}{RGB}{255, 243, 220}

% --- Lists ---
\usepackage{enumitem}

% --- TikZ & pgfplots ---
\usepackage{tikz}
\usepackage{pgfplots}
\pgfplotsset{compat=1.18}
\usetikzlibrary{positioning, arrows.meta, calc, patterns}

% --- Boxes: tcolorbox ---
\usepackage{tcolorbox}
\tcbuselibrary{skins, breakable}

% Generic exercise box (blue border, gray background)
\tcbset{
  exercisebox/.style={
    enhanced,
    colback=accentgray,
    colframe=mainblue,
    boxrule=0.8pt,
    arc=4pt,
    left=10pt, right=10pt, top=8pt, bottom=8pt,
    fonttitle=\bfseries\color{mainblue},
  }
}

% Solution box (green, used in answer sheets)
\tcbset{
  solutionbox/.style={
    enhanced,
    colback=lightgreen,
    colframe=solutiongreen,
    boxrule=0.8pt,
    arc=4pt,
    left=10pt, right=10pt, top=8pt, bottom=8pt,
    fonttitle=\bfseries\color{solutiongreen},
  }
}

% Answer highlight box (orange, for final numeric answers)
\tcbset{
  answerbox/.style={
    enhanced,
    colback=lightorange,
    colframe=warnorange,
    boxrule=0.7pt,
    arc=4pt,
    left=8pt, right=8pt, top=5pt, bottom=5pt,
    fonttitle=\bfseries\color{warnorange},
  }
}

% Step box (light blue border, gray background, for intermediate steps)
\tcbset{
  stepbox/.style={
    enhanced,
    colback=accentgray,
    colframe=mainblue!40,
    boxrule=0.5pt,
    arc=3pt,
    left=8pt, right=8pt, top=6pt, bottom=6pt,
  }
}

% --- Inline note / warning box (mdframed) ---
\usepackage{mdframed}
\newmdenv[
  linecolor=warnorange,
  backgroundcolor=orange!5,
  roundcorner=4pt,
  innerleftmargin=8pt,
  innerrightmargin=8pt,
  innertopmargin=6pt,
  innerbottommargin=6pt,
  linewidth=0.7pt
]{notebox}

% --- Section formatting ---
\usepackage{titlesec}

% Section style: "Exercise N : Title" with a blue rule underneath
% Change the prefix via \renewcommand{\sectionprefix}{...} before \section
\newcommand{\sectionprefix}{Exercise}
\titleformat{\section}[block]
  {\large\bfseries\color{mainblue}}
  {\sectionprefix~\thesection\,:}{0.6em}{}
  [\vspace{-0.3em}\textcolor{mainblue}{\rule{\linewidth}{0.8pt}}]
\titlespacing{\section}{0pt}{1.6em}{0.8em}

% --- Header / Footer ---
\usepackage{fancyhdr}
\pagestyle{fancy}
\fancyhf{}
\renewcommand{\headrulewidth}{0.4pt}

% Customise these three commands in your document after \input{preamble}:
%   \renewcommand{\doccoursename}{My Course}
%   \renewcommand{\docsessionname}{Session N — Topic}
\newcommand{\doccoursename}{Neural Networks for Engineers}
\newcommand{\docsessionname}{Session}

\fancyhead[L]{\small\textcolor{darkgray}{\doccoursename}}
\fancyhead[R]{\small\textcolor{darkgray}{\docsessionname}}
\fancyfoot[C]{\small\textcolor{darkgray}{\thepage}}

% --- Title block macro ---
% Usage: \maketitleblock{Title line}{Subtitle / metadata line}{background color}
% Examples:
%   \maketitleblock{Session 4 — Exercises}{Mixed level | Neural Networks}{mainblue}
%   \maketitleblock{Session 4 — Solutions}{Perceptrons and MLPs | Neural Networks}{solutiongreen}
\newcommand{\maketitleblock}[3]{%
  \begin{tcolorbox}[
    enhanced,
    colback=#3,
    colframe=#3,
    arc=6pt, boxrule=0pt,
    left=16pt, right=16pt, top=12pt, bottom=12pt
  ]
    {\large\bfseries\color{white} #1}\\[4pt]
    {\small\color{white!85!black} #2}
  \end{tcolorbox}
  \vspace{1em}
}

% --- Convenience macros ---

% Blank answer field (inline underline for fill-in tables)
\newcommand{\acell}{\rule{1.8cm}{0.4pt}}

% Step-by-step computation display: \step{label}{math}
% Example: \step{Hidden layer}{z^{(1)} = Wx + b = \ldots}
\newcommand{\step}[2]{%
  \noindent\textbf{#1:}\quad $#2$\par\vspace{0.3em}
}

% Neuron style for TikZ diagrams (reuse with \node[neuron])
% Colors are overridable per layer — see examples below.
%
% Input neuron  : fill=lightblue,  draw=mainblue
% Hidden neuron : fill=green!15,   draw=green!50!black
% Output neuron : fill=orange!15,  draw=orange!70!black
%
% Example architecture diagram:
%
%   \begin{tikzpicture}[
%     neuron/.style={circle, minimum size=1cm, thick, font=\small\bfseries},
%     arr/.style={->, >=Stealth, thick, gray!70}
%   ]
%     \node[neuron, fill=lightblue,  draw=mainblue]          (x1) at (0, 1)  {$x_1$};
%     \node[neuron, fill=green!15,   draw=green!50!black]    (h1) at (3, 1)  {$h_1$};
%     \node[neuron, fill=orange!15,  draw=orange!70!black]   (y)  at (6, 1)  {$\hat{y}$};
%     \draw[arr] (x1) -- (h1);
%     \draw[arr] (h1) -- (y);
%   \end{tikzpicture}
:
%   \renewcommand{\doccoursename}{My Course}
%   \renewcommand{\docsessionname}{Session N — Topic}
\newcommand{\doccoursename}{Neural Networks for Engineers}
\newcommand{\docsessionname}{Session}

\fancyhead[L]{\small\textcolor{darkgray}{\doccoursename}}
\fancyhead[R]{\small\textcolor{darkgray}{\docsessionname}}
\fancyfoot[C]{\small\textcolor{darkgray}{\thepage}}

% --- Title block macro ---
% Usage: \maketitleblock{Title line}{Subtitle / metadata line}{background color}
% Examples:
%   \maketitleblock{Session 4 — Exercises}{Mixed level | Neural Networks}{mainblue}
%   \maketitleblock{Session 4 — Solutions}{Perceptrons and MLPs | Neural Networks}{solutiongreen}
\newcommand{\maketitleblock}[3]{%
  \begin{tcolorbox}[
    enhanced,
    colback=#3,
    colframe=#3,
    arc=6pt, boxrule=0pt,
    left=16pt, right=16pt, top=12pt, bottom=12pt
  ]
    {\large\bfseries\color{white} #1}\\[4pt]
    {\small\color{white!85!black} #2}
  \end{tcolorbox}
  \vspace{1em}
}

% --- Convenience macros ---

% Blank answer field (inline underline for fill-in tables)
\newcommand{\acell}{\rule{1.8cm}{0.4pt}}

% Step-by-step computation display: \step{label}{math}
% Example: \step{Hidden layer}{z^{(1)} = Wx + b = \ldots}
\newcommand{\step}[2]{%
  \noindent\textbf{#1:}\quad $#2$\par\vspace{0.3em}
}

% Neuron style for TikZ diagrams (reuse with \node[neuron])
% Colors are overridable per layer — see examples below.
%
% Input neuron  : fill=lightblue,  draw=mainblue
% Hidden neuron : fill=green!15,   draw=green!50!black
% Output neuron : fill=orange!15,  draw=orange!70!black
%
% Example architecture diagram:
%
%   \begin{tikzpicture}[
%     neuron/.style={circle, minimum size=1cm, thick, font=\small\bfseries},
%     arr/.style={->, >=Stealth, thick, gray!70}
%   ]
%     \node[neuron, fill=lightblue,  draw=mainblue]          (x1) at (0, 1)  {$x_1$};
%     \node[neuron, fill=green!15,   draw=green!50!black]    (h1) at (3, 1)  {$h_1$};
%     \node[neuron, fill=orange!15,  draw=orange!70!black]   (y)  at (6, 1)  {$\hat{y}$};
%     \draw[arr] (x1) -- (h1);
%     \draw[arr] (h1) -- (y);
%   \end{tikzpicture}
 at the top of your document,
%         after \documentclass but before \begin{document}.
% =============================================================

% --- Encoding & language ---
\usepackage[T1]{fontenc}
\usepackage[utf8]{inputenc}
\usepackage[english]{babel}         % swap for [french] if needed

% --- Page geometry ---
\usepackage{geometry}
\geometry{top=2.5cm, bottom=2.5cm, left=2.5cm, right=2.5cm}

% --- Math ---
\usepackage{amsmath, amssymb}

% --- Tables ---
\usepackage{booktabs}               % \toprule, \midrule, \bottomrule
\usepackage{array}                  % column type extensions
\usepackage{tabularx}               % auto-width columns
\usepackage{multirow}               % multi-row cells

% --- Colors ---
\usepackage{xcolor}

\definecolor{mainblue}{RGB}{30, 80, 160}
\definecolor{lightblue}{RGB}{220, 235, 255}
\definecolor{accentgray}{RGB}{245, 246, 248}
\definecolor{darkgray}{RGB}{80, 80, 80}
\definecolor{solutiongreen}{RGB}{20, 120, 60}
\definecolor{lightgreen}{RGB}{220, 245, 230}
\definecolor{warnorange}{RGB}{210, 120, 20}
\definecolor{lightorange}{RGB}{255, 243, 220}

% --- Lists ---
\usepackage{enumitem}

% --- TikZ & pgfplots ---
\usepackage{tikz}
\usepackage{pgfplots}
\pgfplotsset{compat=1.18}
\usetikzlibrary{positioning, arrows.meta, calc, patterns}

% --- Boxes: tcolorbox ---
\usepackage{tcolorbox}
\tcbuselibrary{skins, breakable}

% Generic exercise box (blue border, gray background)
\tcbset{
  exercisebox/.style={
    enhanced,
    colback=accentgray,
    colframe=mainblue,
    boxrule=0.8pt,
    arc=4pt,
    left=10pt, right=10pt, top=8pt, bottom=8pt,
    fonttitle=\bfseries\color{mainblue},
  }
}

% Solution box (green, used in answer sheets)
\tcbset{
  solutionbox/.style={
    enhanced,
    colback=lightgreen,
    colframe=solutiongreen,
    boxrule=0.8pt,
    arc=4pt,
    left=10pt, right=10pt, top=8pt, bottom=8pt,
    fonttitle=\bfseries\color{solutiongreen},
  }
}

% Answer highlight box (orange, for final numeric answers)
\tcbset{
  answerbox/.style={
    enhanced,
    colback=lightorange,
    colframe=warnorange,
    boxrule=0.7pt,
    arc=4pt,
    left=8pt, right=8pt, top=5pt, bottom=5pt,
    fonttitle=\bfseries\color{warnorange},
  }
}

% Step box (light blue border, gray background, for intermediate steps)
\tcbset{
  stepbox/.style={
    enhanced,
    colback=accentgray,
    colframe=mainblue!40,
    boxrule=0.5pt,
    arc=3pt,
    left=8pt, right=8pt, top=6pt, bottom=6pt,
  }
}

% --- Inline note / warning box (mdframed) ---
\usepackage{mdframed}
\newmdenv[
  linecolor=warnorange,
  backgroundcolor=orange!5,
  roundcorner=4pt,
  innerleftmargin=8pt,
  innerrightmargin=8pt,
  innertopmargin=6pt,
  innerbottommargin=6pt,
  linewidth=0.7pt
]{notebox}

% --- Section formatting ---
\usepackage{titlesec}

% Section style: "Exercise N : Title" with a blue rule underneath
% Change the prefix via \renewcommand{\sectionprefix}{...} before \section
\newcommand{\sectionprefix}{Exercise}
\titleformat{\section}[block]
  {\large\bfseries\color{mainblue}}
  {\sectionprefix~\thesection\,:}{0.6em}{}
  [\vspace{-0.3em}\textcolor{mainblue}{\rule{\linewidth}{0.8pt}}]
\titlespacing{\section}{0pt}{1.6em}{0.8em}

% --- Header / Footer ---
\usepackage{fancyhdr}
\pagestyle{fancy}
\fancyhf{}
\renewcommand{\headrulewidth}{0.4pt}

% Customise these three commands in your document after % =============================================================
%  preamble.tex — Reusable preamble for Neural Networks course
%  Usage: % =============================================================
%  preamble.tex — Reusable preamble for Neural Networks course
%  Usage: \input{preamble} at the top of your document,
%         after \documentclass but before \begin{document}.
% =============================================================

% --- Encoding & language ---
\usepackage[T1]{fontenc}
\usepackage[utf8]{inputenc}
\usepackage[english]{babel}         % swap for [french] if needed

% --- Page geometry ---
\usepackage{geometry}
\geometry{top=2.5cm, bottom=2.5cm, left=2.5cm, right=2.5cm}

% --- Math ---
\usepackage{amsmath, amssymb}

% --- Tables ---
\usepackage{booktabs}               % \toprule, \midrule, \bottomrule
\usepackage{array}                  % column type extensions
\usepackage{tabularx}               % auto-width columns
\usepackage{multirow}               % multi-row cells

% --- Colors ---
\usepackage{xcolor}

\definecolor{mainblue}{RGB}{30, 80, 160}
\definecolor{lightblue}{RGB}{220, 235, 255}
\definecolor{accentgray}{RGB}{245, 246, 248}
\definecolor{darkgray}{RGB}{80, 80, 80}
\definecolor{solutiongreen}{RGB}{20, 120, 60}
\definecolor{lightgreen}{RGB}{220, 245, 230}
\definecolor{warnorange}{RGB}{210, 120, 20}
\definecolor{lightorange}{RGB}{255, 243, 220}

% --- Lists ---
\usepackage{enumitem}

% --- TikZ & pgfplots ---
\usepackage{tikz}
\usepackage{pgfplots}
\pgfplotsset{compat=1.18}
\usetikzlibrary{positioning, arrows.meta, calc, patterns}

% --- Boxes: tcolorbox ---
\usepackage{tcolorbox}
\tcbuselibrary{skins, breakable}

% Generic exercise box (blue border, gray background)
\tcbset{
  exercisebox/.style={
    enhanced,
    colback=accentgray,
    colframe=mainblue,
    boxrule=0.8pt,
    arc=4pt,
    left=10pt, right=10pt, top=8pt, bottom=8pt,
    fonttitle=\bfseries\color{mainblue},
  }
}

% Solution box (green, used in answer sheets)
\tcbset{
  solutionbox/.style={
    enhanced,
    colback=lightgreen,
    colframe=solutiongreen,
    boxrule=0.8pt,
    arc=4pt,
    left=10pt, right=10pt, top=8pt, bottom=8pt,
    fonttitle=\bfseries\color{solutiongreen},
  }
}

% Answer highlight box (orange, for final numeric answers)
\tcbset{
  answerbox/.style={
    enhanced,
    colback=lightorange,
    colframe=warnorange,
    boxrule=0.7pt,
    arc=4pt,
    left=8pt, right=8pt, top=5pt, bottom=5pt,
    fonttitle=\bfseries\color{warnorange},
  }
}

% Step box (light blue border, gray background, for intermediate steps)
\tcbset{
  stepbox/.style={
    enhanced,
    colback=accentgray,
    colframe=mainblue!40,
    boxrule=0.5pt,
    arc=3pt,
    left=8pt, right=8pt, top=6pt, bottom=6pt,
  }
}

% --- Inline note / warning box (mdframed) ---
\usepackage{mdframed}
\newmdenv[
  linecolor=warnorange,
  backgroundcolor=orange!5,
  roundcorner=4pt,
  innerleftmargin=8pt,
  innerrightmargin=8pt,
  innertopmargin=6pt,
  innerbottommargin=6pt,
  linewidth=0.7pt
]{notebox}

% --- Section formatting ---
\usepackage{titlesec}

% Section style: "Exercise N : Title" with a blue rule underneath
% Change the prefix via \renewcommand{\sectionprefix}{...} before \section
\newcommand{\sectionprefix}{Exercise}
\titleformat{\section}[block]
  {\large\bfseries\color{mainblue}}
  {\sectionprefix~\thesection\,:}{0.6em}{}
  [\vspace{-0.3em}\textcolor{mainblue}{\rule{\linewidth}{0.8pt}}]
\titlespacing{\section}{0pt}{1.6em}{0.8em}

% --- Header / Footer ---
\usepackage{fancyhdr}
\pagestyle{fancy}
\fancyhf{}
\renewcommand{\headrulewidth}{0.4pt}

% Customise these three commands in your document after \input{preamble}:
%   \renewcommand{\doccoursename}{My Course}
%   \renewcommand{\docsessionname}{Session N — Topic}
\newcommand{\doccoursename}{Neural Networks for Engineers}
\newcommand{\docsessionname}{Session}

\fancyhead[L]{\small\textcolor{darkgray}{\doccoursename}}
\fancyhead[R]{\small\textcolor{darkgray}{\docsessionname}}
\fancyfoot[C]{\small\textcolor{darkgray}{\thepage}}

% --- Title block macro ---
% Usage: \maketitleblock{Title line}{Subtitle / metadata line}{background color}
% Examples:
%   \maketitleblock{Session 4 — Exercises}{Mixed level | Neural Networks}{mainblue}
%   \maketitleblock{Session 4 — Solutions}{Perceptrons and MLPs | Neural Networks}{solutiongreen}
\newcommand{\maketitleblock}[3]{%
  \begin{tcolorbox}[
    enhanced,
    colback=#3,
    colframe=#3,
    arc=6pt, boxrule=0pt,
    left=16pt, right=16pt, top=12pt, bottom=12pt
  ]
    {\large\bfseries\color{white} #1}\\[4pt]
    {\small\color{white!85!black} #2}
  \end{tcolorbox}
  \vspace{1em}
}

% --- Convenience macros ---

% Blank answer field (inline underline for fill-in tables)
\newcommand{\acell}{\rule{1.8cm}{0.4pt}}

% Step-by-step computation display: \step{label}{math}
% Example: \step{Hidden layer}{z^{(1)} = Wx + b = \ldots}
\newcommand{\step}[2]{%
  \noindent\textbf{#1:}\quad $#2$\par\vspace{0.3em}
}

% Neuron style for TikZ diagrams (reuse with \node[neuron])
% Colors are overridable per layer — see examples below.
%
% Input neuron  : fill=lightblue,  draw=mainblue
% Hidden neuron : fill=green!15,   draw=green!50!black
% Output neuron : fill=orange!15,  draw=orange!70!black
%
% Example architecture diagram:
%
%   \begin{tikzpicture}[
%     neuron/.style={circle, minimum size=1cm, thick, font=\small\bfseries},
%     arr/.style={->, >=Stealth, thick, gray!70}
%   ]
%     \node[neuron, fill=lightblue,  draw=mainblue]          (x1) at (0, 1)  {$x_1$};
%     \node[neuron, fill=green!15,   draw=green!50!black]    (h1) at (3, 1)  {$h_1$};
%     \node[neuron, fill=orange!15,  draw=orange!70!black]   (y)  at (6, 1)  {$\hat{y}$};
%     \draw[arr] (x1) -- (h1);
%     \draw[arr] (h1) -- (y);
%   \end{tikzpicture}
 at the top of your document,
%         after \documentclass but before \begin{document}.
% =============================================================

% --- Encoding & language ---
\usepackage[T1]{fontenc}
\usepackage[utf8]{inputenc}
\usepackage[english]{babel}         % swap for [french] if needed

% --- Page geometry ---
\usepackage{geometry}
\geometry{top=2.5cm, bottom=2.5cm, left=2.5cm, right=2.5cm}

% --- Math ---
\usepackage{amsmath, amssymb}

% --- Tables ---
\usepackage{booktabs}               % \toprule, \midrule, \bottomrule
\usepackage{array}                  % column type extensions
\usepackage{tabularx}               % auto-width columns
\usepackage{multirow}               % multi-row cells

% --- Colors ---
\usepackage{xcolor}

\definecolor{mainblue}{RGB}{30, 80, 160}
\definecolor{lightblue}{RGB}{220, 235, 255}
\definecolor{accentgray}{RGB}{245, 246, 248}
\definecolor{darkgray}{RGB}{80, 80, 80}
\definecolor{solutiongreen}{RGB}{20, 120, 60}
\definecolor{lightgreen}{RGB}{220, 245, 230}
\definecolor{warnorange}{RGB}{210, 120, 20}
\definecolor{lightorange}{RGB}{255, 243, 220}

% --- Lists ---
\usepackage{enumitem}

% --- TikZ & pgfplots ---
\usepackage{tikz}
\usepackage{pgfplots}
\pgfplotsset{compat=1.18}
\usetikzlibrary{positioning, arrows.meta, calc, patterns}

% --- Boxes: tcolorbox ---
\usepackage{tcolorbox}
\tcbuselibrary{skins, breakable}

% Generic exercise box (blue border, gray background)
\tcbset{
  exercisebox/.style={
    enhanced,
    colback=accentgray,
    colframe=mainblue,
    boxrule=0.8pt,
    arc=4pt,
    left=10pt, right=10pt, top=8pt, bottom=8pt,
    fonttitle=\bfseries\color{mainblue},
  }
}

% Solution box (green, used in answer sheets)
\tcbset{
  solutionbox/.style={
    enhanced,
    colback=lightgreen,
    colframe=solutiongreen,
    boxrule=0.8pt,
    arc=4pt,
    left=10pt, right=10pt, top=8pt, bottom=8pt,
    fonttitle=\bfseries\color{solutiongreen},
  }
}

% Answer highlight box (orange, for final numeric answers)
\tcbset{
  answerbox/.style={
    enhanced,
    colback=lightorange,
    colframe=warnorange,
    boxrule=0.7pt,
    arc=4pt,
    left=8pt, right=8pt, top=5pt, bottom=5pt,
    fonttitle=\bfseries\color{warnorange},
  }
}

% Step box (light blue border, gray background, for intermediate steps)
\tcbset{
  stepbox/.style={
    enhanced,
    colback=accentgray,
    colframe=mainblue!40,
    boxrule=0.5pt,
    arc=3pt,
    left=8pt, right=8pt, top=6pt, bottom=6pt,
  }
}

% --- Inline note / warning box (mdframed) ---
\usepackage{mdframed}
\newmdenv[
  linecolor=warnorange,
  backgroundcolor=orange!5,
  roundcorner=4pt,
  innerleftmargin=8pt,
  innerrightmargin=8pt,
  innertopmargin=6pt,
  innerbottommargin=6pt,
  linewidth=0.7pt
]{notebox}

% --- Section formatting ---
\usepackage{titlesec}

% Section style: "Exercise N : Title" with a blue rule underneath
% Change the prefix via \renewcommand{\sectionprefix}{...} before \section
\newcommand{\sectionprefix}{Exercise}
\titleformat{\section}[block]
  {\large\bfseries\color{mainblue}}
  {\sectionprefix~\thesection\,:}{0.6em}{}
  [\vspace{-0.3em}\textcolor{mainblue}{\rule{\linewidth}{0.8pt}}]
\titlespacing{\section}{0pt}{1.6em}{0.8em}

% --- Header / Footer ---
\usepackage{fancyhdr}
\pagestyle{fancy}
\fancyhf{}
\renewcommand{\headrulewidth}{0.4pt}

% Customise these three commands in your document after % =============================================================
%  preamble.tex — Reusable preamble for Neural Networks course
%  Usage: \input{preamble} at the top of your document,
%         after \documentclass but before \begin{document}.
% =============================================================

% --- Encoding & language ---
\usepackage[T1]{fontenc}
\usepackage[utf8]{inputenc}
\usepackage[english]{babel}         % swap for [french] if needed

% --- Page geometry ---
\usepackage{geometry}
\geometry{top=2.5cm, bottom=2.5cm, left=2.5cm, right=2.5cm}

% --- Math ---
\usepackage{amsmath, amssymb}

% --- Tables ---
\usepackage{booktabs}               % \toprule, \midrule, \bottomrule
\usepackage{array}                  % column type extensions
\usepackage{tabularx}               % auto-width columns
\usepackage{multirow}               % multi-row cells

% --- Colors ---
\usepackage{xcolor}

\definecolor{mainblue}{RGB}{30, 80, 160}
\definecolor{lightblue}{RGB}{220, 235, 255}
\definecolor{accentgray}{RGB}{245, 246, 248}
\definecolor{darkgray}{RGB}{80, 80, 80}
\definecolor{solutiongreen}{RGB}{20, 120, 60}
\definecolor{lightgreen}{RGB}{220, 245, 230}
\definecolor{warnorange}{RGB}{210, 120, 20}
\definecolor{lightorange}{RGB}{255, 243, 220}

% --- Lists ---
\usepackage{enumitem}

% --- TikZ & pgfplots ---
\usepackage{tikz}
\usepackage{pgfplots}
\pgfplotsset{compat=1.18}
\usetikzlibrary{positioning, arrows.meta, calc, patterns}

% --- Boxes: tcolorbox ---
\usepackage{tcolorbox}
\tcbuselibrary{skins, breakable}

% Generic exercise box (blue border, gray background)
\tcbset{
  exercisebox/.style={
    enhanced,
    colback=accentgray,
    colframe=mainblue,
    boxrule=0.8pt,
    arc=4pt,
    left=10pt, right=10pt, top=8pt, bottom=8pt,
    fonttitle=\bfseries\color{mainblue},
  }
}

% Solution box (green, used in answer sheets)
\tcbset{
  solutionbox/.style={
    enhanced,
    colback=lightgreen,
    colframe=solutiongreen,
    boxrule=0.8pt,
    arc=4pt,
    left=10pt, right=10pt, top=8pt, bottom=8pt,
    fonttitle=\bfseries\color{solutiongreen},
  }
}

% Answer highlight box (orange, for final numeric answers)
\tcbset{
  answerbox/.style={
    enhanced,
    colback=lightorange,
    colframe=warnorange,
    boxrule=0.7pt,
    arc=4pt,
    left=8pt, right=8pt, top=5pt, bottom=5pt,
    fonttitle=\bfseries\color{warnorange},
  }
}

% Step box (light blue border, gray background, for intermediate steps)
\tcbset{
  stepbox/.style={
    enhanced,
    colback=accentgray,
    colframe=mainblue!40,
    boxrule=0.5pt,
    arc=3pt,
    left=8pt, right=8pt, top=6pt, bottom=6pt,
  }
}

% --- Inline note / warning box (mdframed) ---
\usepackage{mdframed}
\newmdenv[
  linecolor=warnorange,
  backgroundcolor=orange!5,
  roundcorner=4pt,
  innerleftmargin=8pt,
  innerrightmargin=8pt,
  innertopmargin=6pt,
  innerbottommargin=6pt,
  linewidth=0.7pt
]{notebox}

% --- Section formatting ---
\usepackage{titlesec}

% Section style: "Exercise N : Title" with a blue rule underneath
% Change the prefix via \renewcommand{\sectionprefix}{...} before \section
\newcommand{\sectionprefix}{Exercise}
\titleformat{\section}[block]
  {\large\bfseries\color{mainblue}}
  {\sectionprefix~\thesection\,:}{0.6em}{}
  [\vspace{-0.3em}\textcolor{mainblue}{\rule{\linewidth}{0.8pt}}]
\titlespacing{\section}{0pt}{1.6em}{0.8em}

% --- Header / Footer ---
\usepackage{fancyhdr}
\pagestyle{fancy}
\fancyhf{}
\renewcommand{\headrulewidth}{0.4pt}

% Customise these three commands in your document after \input{preamble}:
%   \renewcommand{\doccoursename}{My Course}
%   \renewcommand{\docsessionname}{Session N — Topic}
\newcommand{\doccoursename}{Neural Networks for Engineers}
\newcommand{\docsessionname}{Session}

\fancyhead[L]{\small\textcolor{darkgray}{\doccoursename}}
\fancyhead[R]{\small\textcolor{darkgray}{\docsessionname}}
\fancyfoot[C]{\small\textcolor{darkgray}{\thepage}}

% --- Title block macro ---
% Usage: \maketitleblock{Title line}{Subtitle / metadata line}{background color}
% Examples:
%   \maketitleblock{Session 4 — Exercises}{Mixed level | Neural Networks}{mainblue}
%   \maketitleblock{Session 4 — Solutions}{Perceptrons and MLPs | Neural Networks}{solutiongreen}
\newcommand{\maketitleblock}[3]{%
  \begin{tcolorbox}[
    enhanced,
    colback=#3,
    colframe=#3,
    arc=6pt, boxrule=0pt,
    left=16pt, right=16pt, top=12pt, bottom=12pt
  ]
    {\large\bfseries\color{white} #1}\\[4pt]
    {\small\color{white!85!black} #2}
  \end{tcolorbox}
  \vspace{1em}
}

% --- Convenience macros ---

% Blank answer field (inline underline for fill-in tables)
\newcommand{\acell}{\rule{1.8cm}{0.4pt}}

% Step-by-step computation display: \step{label}{math}
% Example: \step{Hidden layer}{z^{(1)} = Wx + b = \ldots}
\newcommand{\step}[2]{%
  \noindent\textbf{#1:}\quad $#2$\par\vspace{0.3em}
}

% Neuron style for TikZ diagrams (reuse with \node[neuron])
% Colors are overridable per layer — see examples below.
%
% Input neuron  : fill=lightblue,  draw=mainblue
% Hidden neuron : fill=green!15,   draw=green!50!black
% Output neuron : fill=orange!15,  draw=orange!70!black
%
% Example architecture diagram:
%
%   \begin{tikzpicture}[
%     neuron/.style={circle, minimum size=1cm, thick, font=\small\bfseries},
%     arr/.style={->, >=Stealth, thick, gray!70}
%   ]
%     \node[neuron, fill=lightblue,  draw=mainblue]          (x1) at (0, 1)  {$x_1$};
%     \node[neuron, fill=green!15,   draw=green!50!black]    (h1) at (3, 1)  {$h_1$};
%     \node[neuron, fill=orange!15,  draw=orange!70!black]   (y)  at (6, 1)  {$\hat{y}$};
%     \draw[arr] (x1) -- (h1);
%     \draw[arr] (h1) -- (y);
%   \end{tikzpicture}
:
%   \renewcommand{\doccoursename}{My Course}
%   \renewcommand{\docsessionname}{Session N — Topic}
\newcommand{\doccoursename}{Neural Networks for Engineers}
\newcommand{\docsessionname}{Session}

\fancyhead[L]{\small\textcolor{darkgray}{\doccoursename}}
\fancyhead[R]{\small\textcolor{darkgray}{\docsessionname}}
\fancyfoot[C]{\small\textcolor{darkgray}{\thepage}}

% --- Title block macro ---
% Usage: \maketitleblock{Title line}{Subtitle / metadata line}{background color}
% Examples:
%   \maketitleblock{Session 4 — Exercises}{Mixed level | Neural Networks}{mainblue}
%   \maketitleblock{Session 4 — Solutions}{Perceptrons and MLPs | Neural Networks}{solutiongreen}
\newcommand{\maketitleblock}[3]{%
  \begin{tcolorbox}[
    enhanced,
    colback=#3,
    colframe=#3,
    arc=6pt, boxrule=0pt,
    left=16pt, right=16pt, top=12pt, bottom=12pt
  ]
    {\large\bfseries\color{white} #1}\\[4pt]
    {\small\color{white!85!black} #2}
  \end{tcolorbox}
  \vspace{1em}
}

% --- Convenience macros ---

% Blank answer field (inline underline for fill-in tables)
\newcommand{\acell}{\rule{1.8cm}{0.4pt}}

% Step-by-step computation display: \step{label}{math}
% Example: \step{Hidden layer}{z^{(1)} = Wx + b = \ldots}
\newcommand{\step}[2]{%
  \noindent\textbf{#1:}\quad $#2$\par\vspace{0.3em}
}

% Neuron style for TikZ diagrams (reuse with \node[neuron])
% Colors are overridable per layer — see examples below.
%
% Input neuron  : fill=lightblue,  draw=mainblue
% Hidden neuron : fill=green!15,   draw=green!50!black
% Output neuron : fill=orange!15,  draw=orange!70!black
%
% Example architecture diagram:
%
%   \begin{tikzpicture}[
%     neuron/.style={circle, minimum size=1cm, thick, font=\small\bfseries},
%     arr/.style={->, >=Stealth, thick, gray!70}
%   ]
%     \node[neuron, fill=lightblue,  draw=mainblue]          (x1) at (0, 1)  {$x_1$};
%     \node[neuron, fill=green!15,   draw=green!50!black]    (h1) at (3, 1)  {$h_1$};
%     \node[neuron, fill=orange!15,  draw=orange!70!black]   (y)  at (6, 1)  {$\hat{y}$};
%     \draw[arr] (x1) -- (h1);
%     \draw[arr] (h1) -- (y);
%   \end{tikzpicture}
:
%   \renewcommand{\doccoursename}{My Course}
%   \renewcommand{\docsessionname}{Session N — Topic}
\newcommand{\doccoursename}{Neural Networks for Engineers}
\newcommand{\docsessionname}{Session}

\fancyhead[L]{\small\textcolor{darkgray}{\doccoursename}}
\fancyhead[R]{\small\textcolor{darkgray}{\docsessionname}}
\fancyfoot[C]{\small\textcolor{darkgray}{\thepage}}

% --- Title block macro ---
% Usage: \maketitleblock{Title line}{Subtitle / metadata line}{background color}
% Examples:
%   \maketitleblock{Session 4 — Exercises}{Mixed level | Neural Networks}{mainblue}
%   \maketitleblock{Session 4 — Solutions}{Perceptrons and MLPs | Neural Networks}{solutiongreen}
\newcommand{\maketitleblock}[3]{%
  \begin{tcolorbox}[
    enhanced,
    colback=#3,
    colframe=#3,
    arc=6pt, boxrule=0pt,
    left=16pt, right=16pt, top=12pt, bottom=12pt
  ]
    {\large\bfseries\color{white} #1}\\[4pt]
    {\small\color{white!85!black} #2}
  \end{tcolorbox}
  \vspace{1em}
}

% --- Convenience macros ---

% Blank answer field (inline underline for fill-in tables)
\newcommand{\acell}{\rule{1.8cm}{0.4pt}}

% Step-by-step computation display: \step{label}{math}
% Example: \step{Hidden layer}{z^{(1)} = Wx + b = \ldots}
\newcommand{\step}[2]{%
  \noindent\textbf{#1:}\quad $#2$\par\vspace{0.3em}
}

% Neuron style for TikZ diagrams (reuse with \node[neuron])
% Colors are overridable per layer — see examples below.
%
% Input neuron  : fill=lightblue,  draw=mainblue
% Hidden neuron : fill=green!15,   draw=green!50!black
% Output neuron : fill=orange!15,  draw=orange!70!black
%
% Example architecture diagram:
%
%   \begin{tikzpicture}[
%     neuron/.style={circle, minimum size=1cm, thick, font=\small\bfseries},
%     arr/.style={->, >=Stealth, thick, gray!70}
%   ]
%     \node[neuron, fill=lightblue,  draw=mainblue]          (x1) at (0, 1)  {$x_1$};
%     \node[neuron, fill=green!15,   draw=green!50!black]    (h1) at (3, 1)  {$h_1$};
%     \node[neuron, fill=orange!15,  draw=orange!70!black]   (y)  at (6, 1)  {$\hat{y}$};
%     \draw[arr] (x1) -- (h1);
%     \draw[arr] (h1) -- (y);
%   \end{tikzpicture}
 at the top of your document,
%         after \documentclass but before \begin{document}.
% =============================================================

% --- Encoding & language ---
\usepackage[T1]{fontenc}
\usepackage[utf8]{inputenc}
\usepackage[english]{babel}         % swap for [french] if needed

% --- Page geometry ---
\usepackage{geometry}
\geometry{top=2.5cm, bottom=2.5cm, left=2.5cm, right=2.5cm}

% --- Math ---
\usepackage{amsmath, amssymb}

% --- Tables ---
\usepackage{booktabs}               % \toprule, \midrule, \bottomrule
\usepackage{array}                  % column type extensions
\usepackage{tabularx}               % auto-width columns
\usepackage{multirow}               % multi-row cells

% --- Colors ---
\usepackage{xcolor}

\definecolor{mainblue}{RGB}{30, 80, 160}
\definecolor{lightblue}{RGB}{220, 235, 255}
\definecolor{accentgray}{RGB}{245, 246, 248}
\definecolor{darkgray}{RGB}{80, 80, 80}
\definecolor{solutiongreen}{RGB}{20, 120, 60}
\definecolor{lightgreen}{RGB}{220, 245, 230}
\definecolor{warnorange}{RGB}{210, 120, 20}
\definecolor{lightorange}{RGB}{255, 243, 220}

% --- Lists ---
\usepackage{enumitem}

% --- TikZ & pgfplots ---
\usepackage{tikz}
\usepackage{pgfplots}
\pgfplotsset{compat=1.18}
\usetikzlibrary{positioning, arrows.meta, calc, patterns}

% --- Boxes: tcolorbox ---
\usepackage{tcolorbox}
\tcbuselibrary{skins, breakable}

% Generic exercise box (blue border, gray background)
\tcbset{
  exercisebox/.style={
    enhanced,
    colback=accentgray,
    colframe=mainblue,
    boxrule=0.8pt,
    arc=4pt,
    left=10pt, right=10pt, top=8pt, bottom=8pt,
    fonttitle=\bfseries\color{mainblue},
  }
}

% Solution box (green, used in answer sheets)
\tcbset{
  solutionbox/.style={
    enhanced,
    colback=lightgreen,
    colframe=solutiongreen,
    boxrule=0.8pt,
    arc=4pt,
    left=10pt, right=10pt, top=8pt, bottom=8pt,
    fonttitle=\bfseries\color{solutiongreen},
  }
}

% Answer highlight box (orange, for final numeric answers)
\tcbset{
  answerbox/.style={
    enhanced,
    colback=lightorange,
    colframe=warnorange,
    boxrule=0.7pt,
    arc=4pt,
    left=8pt, right=8pt, top=5pt, bottom=5pt,
    fonttitle=\bfseries\color{warnorange},
  }
}

% Step box (light blue border, gray background, for intermediate steps)
\tcbset{
  stepbox/.style={
    enhanced,
    colback=accentgray,
    colframe=mainblue!40,
    boxrule=0.5pt,
    arc=3pt,
    left=8pt, right=8pt, top=6pt, bottom=6pt,
  }
}

% --- Inline note / warning box (mdframed) ---
\usepackage{mdframed}
\newmdenv[
  linecolor=warnorange,
  backgroundcolor=orange!5,
  roundcorner=4pt,
  innerleftmargin=8pt,
  innerrightmargin=8pt,
  innertopmargin=6pt,
  innerbottommargin=6pt,
  linewidth=0.7pt
]{notebox}

% --- Section formatting ---
\usepackage{titlesec}

% Section style: "Exercise N : Title" with a blue rule underneath
% Change the prefix via \renewcommand{\sectionprefix}{...} before \section
\newcommand{\sectionprefix}{Exercise}
\titleformat{\section}[block]
  {\large\bfseries\color{mainblue}}
  {\sectionprefix~\thesection\,:}{0.6em}{}
  [\vspace{-0.3em}\textcolor{mainblue}{\rule{\linewidth}{0.8pt}}]
\titlespacing{\section}{0pt}{1.6em}{0.8em}

% --- Header / Footer ---
\usepackage{fancyhdr}
\pagestyle{fancy}
\fancyhf{}
\renewcommand{\headrulewidth}{0.4pt}

% Customise these three commands in your document after % =============================================================
%  preamble.tex — Reusable preamble for Neural Networks course
%  Usage: % =============================================================
%  preamble.tex — Reusable preamble for Neural Networks course
%  Usage: % =============================================================
%  preamble.tex — Reusable preamble for Neural Networks course
%  Usage: \input{preamble} at the top of your document,
%         after \documentclass but before \begin{document}.
% =============================================================

% --- Encoding & language ---
\usepackage[T1]{fontenc}
\usepackage[utf8]{inputenc}
\usepackage[english]{babel}         % swap for [french] if needed

% --- Page geometry ---
\usepackage{geometry}
\geometry{top=2.5cm, bottom=2.5cm, left=2.5cm, right=2.5cm}

% --- Math ---
\usepackage{amsmath, amssymb}

% --- Tables ---
\usepackage{booktabs}               % \toprule, \midrule, \bottomrule
\usepackage{array}                  % column type extensions
\usepackage{tabularx}               % auto-width columns
\usepackage{multirow}               % multi-row cells

% --- Colors ---
\usepackage{xcolor}

\definecolor{mainblue}{RGB}{30, 80, 160}
\definecolor{lightblue}{RGB}{220, 235, 255}
\definecolor{accentgray}{RGB}{245, 246, 248}
\definecolor{darkgray}{RGB}{80, 80, 80}
\definecolor{solutiongreen}{RGB}{20, 120, 60}
\definecolor{lightgreen}{RGB}{220, 245, 230}
\definecolor{warnorange}{RGB}{210, 120, 20}
\definecolor{lightorange}{RGB}{255, 243, 220}

% --- Lists ---
\usepackage{enumitem}

% --- TikZ & pgfplots ---
\usepackage{tikz}
\usepackage{pgfplots}
\pgfplotsset{compat=1.18}
\usetikzlibrary{positioning, arrows.meta, calc, patterns}

% --- Boxes: tcolorbox ---
\usepackage{tcolorbox}
\tcbuselibrary{skins, breakable}

% Generic exercise box (blue border, gray background)
\tcbset{
  exercisebox/.style={
    enhanced,
    colback=accentgray,
    colframe=mainblue,
    boxrule=0.8pt,
    arc=4pt,
    left=10pt, right=10pt, top=8pt, bottom=8pt,
    fonttitle=\bfseries\color{mainblue},
  }
}

% Solution box (green, used in answer sheets)
\tcbset{
  solutionbox/.style={
    enhanced,
    colback=lightgreen,
    colframe=solutiongreen,
    boxrule=0.8pt,
    arc=4pt,
    left=10pt, right=10pt, top=8pt, bottom=8pt,
    fonttitle=\bfseries\color{solutiongreen},
  }
}

% Answer highlight box (orange, for final numeric answers)
\tcbset{
  answerbox/.style={
    enhanced,
    colback=lightorange,
    colframe=warnorange,
    boxrule=0.7pt,
    arc=4pt,
    left=8pt, right=8pt, top=5pt, bottom=5pt,
    fonttitle=\bfseries\color{warnorange},
  }
}

% Step box (light blue border, gray background, for intermediate steps)
\tcbset{
  stepbox/.style={
    enhanced,
    colback=accentgray,
    colframe=mainblue!40,
    boxrule=0.5pt,
    arc=3pt,
    left=8pt, right=8pt, top=6pt, bottom=6pt,
  }
}

% --- Inline note / warning box (mdframed) ---
\usepackage{mdframed}
\newmdenv[
  linecolor=warnorange,
  backgroundcolor=orange!5,
  roundcorner=4pt,
  innerleftmargin=8pt,
  innerrightmargin=8pt,
  innertopmargin=6pt,
  innerbottommargin=6pt,
  linewidth=0.7pt
]{notebox}

% --- Section formatting ---
\usepackage{titlesec}

% Section style: "Exercise N : Title" with a blue rule underneath
% Change the prefix via \renewcommand{\sectionprefix}{...} before \section
\newcommand{\sectionprefix}{Exercise}
\titleformat{\section}[block]
  {\large\bfseries\color{mainblue}}
  {\sectionprefix~\thesection\,:}{0.6em}{}
  [\vspace{-0.3em}\textcolor{mainblue}{\rule{\linewidth}{0.8pt}}]
\titlespacing{\section}{0pt}{1.6em}{0.8em}

% --- Header / Footer ---
\usepackage{fancyhdr}
\pagestyle{fancy}
\fancyhf{}
\renewcommand{\headrulewidth}{0.4pt}

% Customise these three commands in your document after \input{preamble}:
%   \renewcommand{\doccoursename}{My Course}
%   \renewcommand{\docsessionname}{Session N — Topic}
\newcommand{\doccoursename}{Neural Networks for Engineers}
\newcommand{\docsessionname}{Session}

\fancyhead[L]{\small\textcolor{darkgray}{\doccoursename}}
\fancyhead[R]{\small\textcolor{darkgray}{\docsessionname}}
\fancyfoot[C]{\small\textcolor{darkgray}{\thepage}}

% --- Title block macro ---
% Usage: \maketitleblock{Title line}{Subtitle / metadata line}{background color}
% Examples:
%   \maketitleblock{Session 4 — Exercises}{Mixed level | Neural Networks}{mainblue}
%   \maketitleblock{Session 4 — Solutions}{Perceptrons and MLPs | Neural Networks}{solutiongreen}
\newcommand{\maketitleblock}[3]{%
  \begin{tcolorbox}[
    enhanced,
    colback=#3,
    colframe=#3,
    arc=6pt, boxrule=0pt,
    left=16pt, right=16pt, top=12pt, bottom=12pt
  ]
    {\large\bfseries\color{white} #1}\\[4pt]
    {\small\color{white!85!black} #2}
  \end{tcolorbox}
  \vspace{1em}
}

% --- Convenience macros ---

% Blank answer field (inline underline for fill-in tables)
\newcommand{\acell}{\rule{1.8cm}{0.4pt}}

% Step-by-step computation display: \step{label}{math}
% Example: \step{Hidden layer}{z^{(1)} = Wx + b = \ldots}
\newcommand{\step}[2]{%
  \noindent\textbf{#1:}\quad $#2$\par\vspace{0.3em}
}

% Neuron style for TikZ diagrams (reuse with \node[neuron])
% Colors are overridable per layer — see examples below.
%
% Input neuron  : fill=lightblue,  draw=mainblue
% Hidden neuron : fill=green!15,   draw=green!50!black
% Output neuron : fill=orange!15,  draw=orange!70!black
%
% Example architecture diagram:
%
%   \begin{tikzpicture}[
%     neuron/.style={circle, minimum size=1cm, thick, font=\small\bfseries},
%     arr/.style={->, >=Stealth, thick, gray!70}
%   ]
%     \node[neuron, fill=lightblue,  draw=mainblue]          (x1) at (0, 1)  {$x_1$};
%     \node[neuron, fill=green!15,   draw=green!50!black]    (h1) at (3, 1)  {$h_1$};
%     \node[neuron, fill=orange!15,  draw=orange!70!black]   (y)  at (6, 1)  {$\hat{y}$};
%     \draw[arr] (x1) -- (h1);
%     \draw[arr] (h1) -- (y);
%   \end{tikzpicture}
 at the top of your document,
%         after \documentclass but before \begin{document}.
% =============================================================

% --- Encoding & language ---
\usepackage[T1]{fontenc}
\usepackage[utf8]{inputenc}
\usepackage[english]{babel}         % swap for [french] if needed

% --- Page geometry ---
\usepackage{geometry}
\geometry{top=2.5cm, bottom=2.5cm, left=2.5cm, right=2.5cm}

% --- Math ---
\usepackage{amsmath, amssymb}

% --- Tables ---
\usepackage{booktabs}               % \toprule, \midrule, \bottomrule
\usepackage{array}                  % column type extensions
\usepackage{tabularx}               % auto-width columns
\usepackage{multirow}               % multi-row cells

% --- Colors ---
\usepackage{xcolor}

\definecolor{mainblue}{RGB}{30, 80, 160}
\definecolor{lightblue}{RGB}{220, 235, 255}
\definecolor{accentgray}{RGB}{245, 246, 248}
\definecolor{darkgray}{RGB}{80, 80, 80}
\definecolor{solutiongreen}{RGB}{20, 120, 60}
\definecolor{lightgreen}{RGB}{220, 245, 230}
\definecolor{warnorange}{RGB}{210, 120, 20}
\definecolor{lightorange}{RGB}{255, 243, 220}

% --- Lists ---
\usepackage{enumitem}

% --- TikZ & pgfplots ---
\usepackage{tikz}
\usepackage{pgfplots}
\pgfplotsset{compat=1.18}
\usetikzlibrary{positioning, arrows.meta, calc, patterns}

% --- Boxes: tcolorbox ---
\usepackage{tcolorbox}
\tcbuselibrary{skins, breakable}

% Generic exercise box (blue border, gray background)
\tcbset{
  exercisebox/.style={
    enhanced,
    colback=accentgray,
    colframe=mainblue,
    boxrule=0.8pt,
    arc=4pt,
    left=10pt, right=10pt, top=8pt, bottom=8pt,
    fonttitle=\bfseries\color{mainblue},
  }
}

% Solution box (green, used in answer sheets)
\tcbset{
  solutionbox/.style={
    enhanced,
    colback=lightgreen,
    colframe=solutiongreen,
    boxrule=0.8pt,
    arc=4pt,
    left=10pt, right=10pt, top=8pt, bottom=8pt,
    fonttitle=\bfseries\color{solutiongreen},
  }
}

% Answer highlight box (orange, for final numeric answers)
\tcbset{
  answerbox/.style={
    enhanced,
    colback=lightorange,
    colframe=warnorange,
    boxrule=0.7pt,
    arc=4pt,
    left=8pt, right=8pt, top=5pt, bottom=5pt,
    fonttitle=\bfseries\color{warnorange},
  }
}

% Step box (light blue border, gray background, for intermediate steps)
\tcbset{
  stepbox/.style={
    enhanced,
    colback=accentgray,
    colframe=mainblue!40,
    boxrule=0.5pt,
    arc=3pt,
    left=8pt, right=8pt, top=6pt, bottom=6pt,
  }
}

% --- Inline note / warning box (mdframed) ---
\usepackage{mdframed}
\newmdenv[
  linecolor=warnorange,
  backgroundcolor=orange!5,
  roundcorner=4pt,
  innerleftmargin=8pt,
  innerrightmargin=8pt,
  innertopmargin=6pt,
  innerbottommargin=6pt,
  linewidth=0.7pt
]{notebox}

% --- Section formatting ---
\usepackage{titlesec}

% Section style: "Exercise N : Title" with a blue rule underneath
% Change the prefix via \renewcommand{\sectionprefix}{...} before \section
\newcommand{\sectionprefix}{Exercise}
\titleformat{\section}[block]
  {\large\bfseries\color{mainblue}}
  {\sectionprefix~\thesection\,:}{0.6em}{}
  [\vspace{-0.3em}\textcolor{mainblue}{\rule{\linewidth}{0.8pt}}]
\titlespacing{\section}{0pt}{1.6em}{0.8em}

% --- Header / Footer ---
\usepackage{fancyhdr}
\pagestyle{fancy}
\fancyhf{}
\renewcommand{\headrulewidth}{0.4pt}

% Customise these three commands in your document after % =============================================================
%  preamble.tex — Reusable preamble for Neural Networks course
%  Usage: \input{preamble} at the top of your document,
%         after \documentclass but before \begin{document}.
% =============================================================

% --- Encoding & language ---
\usepackage[T1]{fontenc}
\usepackage[utf8]{inputenc}
\usepackage[english]{babel}         % swap for [french] if needed

% --- Page geometry ---
\usepackage{geometry}
\geometry{top=2.5cm, bottom=2.5cm, left=2.5cm, right=2.5cm}

% --- Math ---
\usepackage{amsmath, amssymb}

% --- Tables ---
\usepackage{booktabs}               % \toprule, \midrule, \bottomrule
\usepackage{array}                  % column type extensions
\usepackage{tabularx}               % auto-width columns
\usepackage{multirow}               % multi-row cells

% --- Colors ---
\usepackage{xcolor}

\definecolor{mainblue}{RGB}{30, 80, 160}
\definecolor{lightblue}{RGB}{220, 235, 255}
\definecolor{accentgray}{RGB}{245, 246, 248}
\definecolor{darkgray}{RGB}{80, 80, 80}
\definecolor{solutiongreen}{RGB}{20, 120, 60}
\definecolor{lightgreen}{RGB}{220, 245, 230}
\definecolor{warnorange}{RGB}{210, 120, 20}
\definecolor{lightorange}{RGB}{255, 243, 220}

% --- Lists ---
\usepackage{enumitem}

% --- TikZ & pgfplots ---
\usepackage{tikz}
\usepackage{pgfplots}
\pgfplotsset{compat=1.18}
\usetikzlibrary{positioning, arrows.meta, calc, patterns}

% --- Boxes: tcolorbox ---
\usepackage{tcolorbox}
\tcbuselibrary{skins, breakable}

% Generic exercise box (blue border, gray background)
\tcbset{
  exercisebox/.style={
    enhanced,
    colback=accentgray,
    colframe=mainblue,
    boxrule=0.8pt,
    arc=4pt,
    left=10pt, right=10pt, top=8pt, bottom=8pt,
    fonttitle=\bfseries\color{mainblue},
  }
}

% Solution box (green, used in answer sheets)
\tcbset{
  solutionbox/.style={
    enhanced,
    colback=lightgreen,
    colframe=solutiongreen,
    boxrule=0.8pt,
    arc=4pt,
    left=10pt, right=10pt, top=8pt, bottom=8pt,
    fonttitle=\bfseries\color{solutiongreen},
  }
}

% Answer highlight box (orange, for final numeric answers)
\tcbset{
  answerbox/.style={
    enhanced,
    colback=lightorange,
    colframe=warnorange,
    boxrule=0.7pt,
    arc=4pt,
    left=8pt, right=8pt, top=5pt, bottom=5pt,
    fonttitle=\bfseries\color{warnorange},
  }
}

% Step box (light blue border, gray background, for intermediate steps)
\tcbset{
  stepbox/.style={
    enhanced,
    colback=accentgray,
    colframe=mainblue!40,
    boxrule=0.5pt,
    arc=3pt,
    left=8pt, right=8pt, top=6pt, bottom=6pt,
  }
}

% --- Inline note / warning box (mdframed) ---
\usepackage{mdframed}
\newmdenv[
  linecolor=warnorange,
  backgroundcolor=orange!5,
  roundcorner=4pt,
  innerleftmargin=8pt,
  innerrightmargin=8pt,
  innertopmargin=6pt,
  innerbottommargin=6pt,
  linewidth=0.7pt
]{notebox}

% --- Section formatting ---
\usepackage{titlesec}

% Section style: "Exercise N : Title" with a blue rule underneath
% Change the prefix via \renewcommand{\sectionprefix}{...} before \section
\newcommand{\sectionprefix}{Exercise}
\titleformat{\section}[block]
  {\large\bfseries\color{mainblue}}
  {\sectionprefix~\thesection\,:}{0.6em}{}
  [\vspace{-0.3em}\textcolor{mainblue}{\rule{\linewidth}{0.8pt}}]
\titlespacing{\section}{0pt}{1.6em}{0.8em}

% --- Header / Footer ---
\usepackage{fancyhdr}
\pagestyle{fancy}
\fancyhf{}
\renewcommand{\headrulewidth}{0.4pt}

% Customise these three commands in your document after \input{preamble}:
%   \renewcommand{\doccoursename}{My Course}
%   \renewcommand{\docsessionname}{Session N — Topic}
\newcommand{\doccoursename}{Neural Networks for Engineers}
\newcommand{\docsessionname}{Session}

\fancyhead[L]{\small\textcolor{darkgray}{\doccoursename}}
\fancyhead[R]{\small\textcolor{darkgray}{\docsessionname}}
\fancyfoot[C]{\small\textcolor{darkgray}{\thepage}}

% --- Title block macro ---
% Usage: \maketitleblock{Title line}{Subtitle / metadata line}{background color}
% Examples:
%   \maketitleblock{Session 4 — Exercises}{Mixed level | Neural Networks}{mainblue}
%   \maketitleblock{Session 4 — Solutions}{Perceptrons and MLPs | Neural Networks}{solutiongreen}
\newcommand{\maketitleblock}[3]{%
  \begin{tcolorbox}[
    enhanced,
    colback=#3,
    colframe=#3,
    arc=6pt, boxrule=0pt,
    left=16pt, right=16pt, top=12pt, bottom=12pt
  ]
    {\large\bfseries\color{white} #1}\\[4pt]
    {\small\color{white!85!black} #2}
  \end{tcolorbox}
  \vspace{1em}
}

% --- Convenience macros ---

% Blank answer field (inline underline for fill-in tables)
\newcommand{\acell}{\rule{1.8cm}{0.4pt}}

% Step-by-step computation display: \step{label}{math}
% Example: \step{Hidden layer}{z^{(1)} = Wx + b = \ldots}
\newcommand{\step}[2]{%
  \noindent\textbf{#1:}\quad $#2$\par\vspace{0.3em}
}

% Neuron style for TikZ diagrams (reuse with \node[neuron])
% Colors are overridable per layer — see examples below.
%
% Input neuron  : fill=lightblue,  draw=mainblue
% Hidden neuron : fill=green!15,   draw=green!50!black
% Output neuron : fill=orange!15,  draw=orange!70!black
%
% Example architecture diagram:
%
%   \begin{tikzpicture}[
%     neuron/.style={circle, minimum size=1cm, thick, font=\small\bfseries},
%     arr/.style={->, >=Stealth, thick, gray!70}
%   ]
%     \node[neuron, fill=lightblue,  draw=mainblue]          (x1) at (0, 1)  {$x_1$};
%     \node[neuron, fill=green!15,   draw=green!50!black]    (h1) at (3, 1)  {$h_1$};
%     \node[neuron, fill=orange!15,  draw=orange!70!black]   (y)  at (6, 1)  {$\hat{y}$};
%     \draw[arr] (x1) -- (h1);
%     \draw[arr] (h1) -- (y);
%   \end{tikzpicture}
:
%   \renewcommand{\doccoursename}{My Course}
%   \renewcommand{\docsessionname}{Session N — Topic}
\newcommand{\doccoursename}{Neural Networks for Engineers}
\newcommand{\docsessionname}{Session}

\fancyhead[L]{\small\textcolor{darkgray}{\doccoursename}}
\fancyhead[R]{\small\textcolor{darkgray}{\docsessionname}}
\fancyfoot[C]{\small\textcolor{darkgray}{\thepage}}

% --- Title block macro ---
% Usage: \maketitleblock{Title line}{Subtitle / metadata line}{background color}
% Examples:
%   \maketitleblock{Session 4 — Exercises}{Mixed level | Neural Networks}{mainblue}
%   \maketitleblock{Session 4 — Solutions}{Perceptrons and MLPs | Neural Networks}{solutiongreen}
\newcommand{\maketitleblock}[3]{%
  \begin{tcolorbox}[
    enhanced,
    colback=#3,
    colframe=#3,
    arc=6pt, boxrule=0pt,
    left=16pt, right=16pt, top=12pt, bottom=12pt
  ]
    {\large\bfseries\color{white} #1}\\[4pt]
    {\small\color{white!85!black} #2}
  \end{tcolorbox}
  \vspace{1em}
}

% --- Convenience macros ---

% Blank answer field (inline underline for fill-in tables)
\newcommand{\acell}{\rule{1.8cm}{0.4pt}}

% Step-by-step computation display: \step{label}{math}
% Example: \step{Hidden layer}{z^{(1)} = Wx + b = \ldots}
\newcommand{\step}[2]{%
  \noindent\textbf{#1:}\quad $#2$\par\vspace{0.3em}
}

% Neuron style for TikZ diagrams (reuse with \node[neuron])
% Colors are overridable per layer — see examples below.
%
% Input neuron  : fill=lightblue,  draw=mainblue
% Hidden neuron : fill=green!15,   draw=green!50!black
% Output neuron : fill=orange!15,  draw=orange!70!black
%
% Example architecture diagram:
%
%   \begin{tikzpicture}[
%     neuron/.style={circle, minimum size=1cm, thick, font=\small\bfseries},
%     arr/.style={->, >=Stealth, thick, gray!70}
%   ]
%     \node[neuron, fill=lightblue,  draw=mainblue]          (x1) at (0, 1)  {$x_1$};
%     \node[neuron, fill=green!15,   draw=green!50!black]    (h1) at (3, 1)  {$h_1$};
%     \node[neuron, fill=orange!15,  draw=orange!70!black]   (y)  at (6, 1)  {$\hat{y}$};
%     \draw[arr] (x1) -- (h1);
%     \draw[arr] (h1) -- (y);
%   \end{tikzpicture}
 at the top of your document,
%         after \documentclass but before \begin{document}.
% =============================================================

% --- Encoding & language ---
\usepackage[T1]{fontenc}
\usepackage[utf8]{inputenc}
\usepackage[english]{babel}         % swap for [french] if needed

% --- Page geometry ---
\usepackage{geometry}
\geometry{top=2.5cm, bottom=2.5cm, left=2.5cm, right=2.5cm}

% --- Math ---
\usepackage{amsmath, amssymb}

% --- Tables ---
\usepackage{booktabs}               % \toprule, \midrule, \bottomrule
\usepackage{array}                  % column type extensions
\usepackage{tabularx}               % auto-width columns
\usepackage{multirow}               % multi-row cells

% --- Colors ---
\usepackage{xcolor}

\definecolor{mainblue}{RGB}{30, 80, 160}
\definecolor{lightblue}{RGB}{220, 235, 255}
\definecolor{accentgray}{RGB}{245, 246, 248}
\definecolor{darkgray}{RGB}{80, 80, 80}
\definecolor{solutiongreen}{RGB}{20, 120, 60}
\definecolor{lightgreen}{RGB}{220, 245, 230}
\definecolor{warnorange}{RGB}{210, 120, 20}
\definecolor{lightorange}{RGB}{255, 243, 220}

% --- Lists ---
\usepackage{enumitem}

% --- TikZ & pgfplots ---
\usepackage{tikz}
\usepackage{pgfplots}
\pgfplotsset{compat=1.18}
\usetikzlibrary{positioning, arrows.meta, calc, patterns}

% --- Boxes: tcolorbox ---
\usepackage{tcolorbox}
\tcbuselibrary{skins, breakable}

% Generic exercise box (blue border, gray background)
\tcbset{
  exercisebox/.style={
    enhanced,
    colback=accentgray,
    colframe=mainblue,
    boxrule=0.8pt,
    arc=4pt,
    left=10pt, right=10pt, top=8pt, bottom=8pt,
    fonttitle=\bfseries\color{mainblue},
  }
}

% Solution box (green, used in answer sheets)
\tcbset{
  solutionbox/.style={
    enhanced,
    colback=lightgreen,
    colframe=solutiongreen,
    boxrule=0.8pt,
    arc=4pt,
    left=10pt, right=10pt, top=8pt, bottom=8pt,
    fonttitle=\bfseries\color{solutiongreen},
  }
}

% Answer highlight box (orange, for final numeric answers)
\tcbset{
  answerbox/.style={
    enhanced,
    colback=lightorange,
    colframe=warnorange,
    boxrule=0.7pt,
    arc=4pt,
    left=8pt, right=8pt, top=5pt, bottom=5pt,
    fonttitle=\bfseries\color{warnorange},
  }
}

% Step box (light blue border, gray background, for intermediate steps)
\tcbset{
  stepbox/.style={
    enhanced,
    colback=accentgray,
    colframe=mainblue!40,
    boxrule=0.5pt,
    arc=3pt,
    left=8pt, right=8pt, top=6pt, bottom=6pt,
  }
}

% --- Inline note / warning box (mdframed) ---
\usepackage{mdframed}
\newmdenv[
  linecolor=warnorange,
  backgroundcolor=orange!5,
  roundcorner=4pt,
  innerleftmargin=8pt,
  innerrightmargin=8pt,
  innertopmargin=6pt,
  innerbottommargin=6pt,
  linewidth=0.7pt
]{notebox}

% --- Section formatting ---
\usepackage{titlesec}

% Section style: "Exercise N : Title" with a blue rule underneath
% Change the prefix via \renewcommand{\sectionprefix}{...} before \section
\newcommand{\sectionprefix}{Exercise}
\titleformat{\section}[block]
  {\large\bfseries\color{mainblue}}
  {\sectionprefix~\thesection\,:}{0.6em}{}
  [\vspace{-0.3em}\textcolor{mainblue}{\rule{\linewidth}{0.8pt}}]
\titlespacing{\section}{0pt}{1.6em}{0.8em}

% --- Header / Footer ---
\usepackage{fancyhdr}
\pagestyle{fancy}
\fancyhf{}
\renewcommand{\headrulewidth}{0.4pt}

% Customise these three commands in your document after % =============================================================
%  preamble.tex — Reusable preamble for Neural Networks course
%  Usage: % =============================================================
%  preamble.tex — Reusable preamble for Neural Networks course
%  Usage: \input{preamble} at the top of your document,
%         after \documentclass but before \begin{document}.
% =============================================================

% --- Encoding & language ---
\usepackage[T1]{fontenc}
\usepackage[utf8]{inputenc}
\usepackage[english]{babel}         % swap for [french] if needed

% --- Page geometry ---
\usepackage{geometry}
\geometry{top=2.5cm, bottom=2.5cm, left=2.5cm, right=2.5cm}

% --- Math ---
\usepackage{amsmath, amssymb}

% --- Tables ---
\usepackage{booktabs}               % \toprule, \midrule, \bottomrule
\usepackage{array}                  % column type extensions
\usepackage{tabularx}               % auto-width columns
\usepackage{multirow}               % multi-row cells

% --- Colors ---
\usepackage{xcolor}

\definecolor{mainblue}{RGB}{30, 80, 160}
\definecolor{lightblue}{RGB}{220, 235, 255}
\definecolor{accentgray}{RGB}{245, 246, 248}
\definecolor{darkgray}{RGB}{80, 80, 80}
\definecolor{solutiongreen}{RGB}{20, 120, 60}
\definecolor{lightgreen}{RGB}{220, 245, 230}
\definecolor{warnorange}{RGB}{210, 120, 20}
\definecolor{lightorange}{RGB}{255, 243, 220}

% --- Lists ---
\usepackage{enumitem}

% --- TikZ & pgfplots ---
\usepackage{tikz}
\usepackage{pgfplots}
\pgfplotsset{compat=1.18}
\usetikzlibrary{positioning, arrows.meta, calc, patterns}

% --- Boxes: tcolorbox ---
\usepackage{tcolorbox}
\tcbuselibrary{skins, breakable}

% Generic exercise box (blue border, gray background)
\tcbset{
  exercisebox/.style={
    enhanced,
    colback=accentgray,
    colframe=mainblue,
    boxrule=0.8pt,
    arc=4pt,
    left=10pt, right=10pt, top=8pt, bottom=8pt,
    fonttitle=\bfseries\color{mainblue},
  }
}

% Solution box (green, used in answer sheets)
\tcbset{
  solutionbox/.style={
    enhanced,
    colback=lightgreen,
    colframe=solutiongreen,
    boxrule=0.8pt,
    arc=4pt,
    left=10pt, right=10pt, top=8pt, bottom=8pt,
    fonttitle=\bfseries\color{solutiongreen},
  }
}

% Answer highlight box (orange, for final numeric answers)
\tcbset{
  answerbox/.style={
    enhanced,
    colback=lightorange,
    colframe=warnorange,
    boxrule=0.7pt,
    arc=4pt,
    left=8pt, right=8pt, top=5pt, bottom=5pt,
    fonttitle=\bfseries\color{warnorange},
  }
}

% Step box (light blue border, gray background, for intermediate steps)
\tcbset{
  stepbox/.style={
    enhanced,
    colback=accentgray,
    colframe=mainblue!40,
    boxrule=0.5pt,
    arc=3pt,
    left=8pt, right=8pt, top=6pt, bottom=6pt,
  }
}

% --- Inline note / warning box (mdframed) ---
\usepackage{mdframed}
\newmdenv[
  linecolor=warnorange,
  backgroundcolor=orange!5,
  roundcorner=4pt,
  innerleftmargin=8pt,
  innerrightmargin=8pt,
  innertopmargin=6pt,
  innerbottommargin=6pt,
  linewidth=0.7pt
]{notebox}

% --- Section formatting ---
\usepackage{titlesec}

% Section style: "Exercise N : Title" with a blue rule underneath
% Change the prefix via \renewcommand{\sectionprefix}{...} before \section
\newcommand{\sectionprefix}{Exercise}
\titleformat{\section}[block]
  {\large\bfseries\color{mainblue}}
  {\sectionprefix~\thesection\,:}{0.6em}{}
  [\vspace{-0.3em}\textcolor{mainblue}{\rule{\linewidth}{0.8pt}}]
\titlespacing{\section}{0pt}{1.6em}{0.8em}

% --- Header / Footer ---
\usepackage{fancyhdr}
\pagestyle{fancy}
\fancyhf{}
\renewcommand{\headrulewidth}{0.4pt}

% Customise these three commands in your document after \input{preamble}:
%   \renewcommand{\doccoursename}{My Course}
%   \renewcommand{\docsessionname}{Session N — Topic}
\newcommand{\doccoursename}{Neural Networks for Engineers}
\newcommand{\docsessionname}{Session}

\fancyhead[L]{\small\textcolor{darkgray}{\doccoursename}}
\fancyhead[R]{\small\textcolor{darkgray}{\docsessionname}}
\fancyfoot[C]{\small\textcolor{darkgray}{\thepage}}

% --- Title block macro ---
% Usage: \maketitleblock{Title line}{Subtitle / metadata line}{background color}
% Examples:
%   \maketitleblock{Session 4 — Exercises}{Mixed level | Neural Networks}{mainblue}
%   \maketitleblock{Session 4 — Solutions}{Perceptrons and MLPs | Neural Networks}{solutiongreen}
\newcommand{\maketitleblock}[3]{%
  \begin{tcolorbox}[
    enhanced,
    colback=#3,
    colframe=#3,
    arc=6pt, boxrule=0pt,
    left=16pt, right=16pt, top=12pt, bottom=12pt
  ]
    {\large\bfseries\color{white} #1}\\[4pt]
    {\small\color{white!85!black} #2}
  \end{tcolorbox}
  \vspace{1em}
}

% --- Convenience macros ---

% Blank answer field (inline underline for fill-in tables)
\newcommand{\acell}{\rule{1.8cm}{0.4pt}}

% Step-by-step computation display: \step{label}{math}
% Example: \step{Hidden layer}{z^{(1)} = Wx + b = \ldots}
\newcommand{\step}[2]{%
  \noindent\textbf{#1:}\quad $#2$\par\vspace{0.3em}
}

% Neuron style for TikZ diagrams (reuse with \node[neuron])
% Colors are overridable per layer — see examples below.
%
% Input neuron  : fill=lightblue,  draw=mainblue
% Hidden neuron : fill=green!15,   draw=green!50!black
% Output neuron : fill=orange!15,  draw=orange!70!black
%
% Example architecture diagram:
%
%   \begin{tikzpicture}[
%     neuron/.style={circle, minimum size=1cm, thick, font=\small\bfseries},
%     arr/.style={->, >=Stealth, thick, gray!70}
%   ]
%     \node[neuron, fill=lightblue,  draw=mainblue]          (x1) at (0, 1)  {$x_1$};
%     \node[neuron, fill=green!15,   draw=green!50!black]    (h1) at (3, 1)  {$h_1$};
%     \node[neuron, fill=orange!15,  draw=orange!70!black]   (y)  at (6, 1)  {$\hat{y}$};
%     \draw[arr] (x1) -- (h1);
%     \draw[arr] (h1) -- (y);
%   \end{tikzpicture}
 at the top of your document,
%         after \documentclass but before \begin{document}.
% =============================================================

% --- Encoding & language ---
\usepackage[T1]{fontenc}
\usepackage[utf8]{inputenc}
\usepackage[english]{babel}         % swap for [french] if needed

% --- Page geometry ---
\usepackage{geometry}
\geometry{top=2.5cm, bottom=2.5cm, left=2.5cm, right=2.5cm}

% --- Math ---
\usepackage{amsmath, amssymb}

% --- Tables ---
\usepackage{booktabs}               % \toprule, \midrule, \bottomrule
\usepackage{array}                  % column type extensions
\usepackage{tabularx}               % auto-width columns
\usepackage{multirow}               % multi-row cells

% --- Colors ---
\usepackage{xcolor}

\definecolor{mainblue}{RGB}{30, 80, 160}
\definecolor{lightblue}{RGB}{220, 235, 255}
\definecolor{accentgray}{RGB}{245, 246, 248}
\definecolor{darkgray}{RGB}{80, 80, 80}
\definecolor{solutiongreen}{RGB}{20, 120, 60}
\definecolor{lightgreen}{RGB}{220, 245, 230}
\definecolor{warnorange}{RGB}{210, 120, 20}
\definecolor{lightorange}{RGB}{255, 243, 220}

% --- Lists ---
\usepackage{enumitem}

% --- TikZ & pgfplots ---
\usepackage{tikz}
\usepackage{pgfplots}
\pgfplotsset{compat=1.18}
\usetikzlibrary{positioning, arrows.meta, calc, patterns}

% --- Boxes: tcolorbox ---
\usepackage{tcolorbox}
\tcbuselibrary{skins, breakable}

% Generic exercise box (blue border, gray background)
\tcbset{
  exercisebox/.style={
    enhanced,
    colback=accentgray,
    colframe=mainblue,
    boxrule=0.8pt,
    arc=4pt,
    left=10pt, right=10pt, top=8pt, bottom=8pt,
    fonttitle=\bfseries\color{mainblue},
  }
}

% Solution box (green, used in answer sheets)
\tcbset{
  solutionbox/.style={
    enhanced,
    colback=lightgreen,
    colframe=solutiongreen,
    boxrule=0.8pt,
    arc=4pt,
    left=10pt, right=10pt, top=8pt, bottom=8pt,
    fonttitle=\bfseries\color{solutiongreen},
  }
}

% Answer highlight box (orange, for final numeric answers)
\tcbset{
  answerbox/.style={
    enhanced,
    colback=lightorange,
    colframe=warnorange,
    boxrule=0.7pt,
    arc=4pt,
    left=8pt, right=8pt, top=5pt, bottom=5pt,
    fonttitle=\bfseries\color{warnorange},
  }
}

% Step box (light blue border, gray background, for intermediate steps)
\tcbset{
  stepbox/.style={
    enhanced,
    colback=accentgray,
    colframe=mainblue!40,
    boxrule=0.5pt,
    arc=3pt,
    left=8pt, right=8pt, top=6pt, bottom=6pt,
  }
}

% --- Inline note / warning box (mdframed) ---
\usepackage{mdframed}
\newmdenv[
  linecolor=warnorange,
  backgroundcolor=orange!5,
  roundcorner=4pt,
  innerleftmargin=8pt,
  innerrightmargin=8pt,
  innertopmargin=6pt,
  innerbottommargin=6pt,
  linewidth=0.7pt
]{notebox}

% --- Section formatting ---
\usepackage{titlesec}

% Section style: "Exercise N : Title" with a blue rule underneath
% Change the prefix via \renewcommand{\sectionprefix}{...} before \section
\newcommand{\sectionprefix}{Exercise}
\titleformat{\section}[block]
  {\large\bfseries\color{mainblue}}
  {\sectionprefix~\thesection\,:}{0.6em}{}
  [\vspace{-0.3em}\textcolor{mainblue}{\rule{\linewidth}{0.8pt}}]
\titlespacing{\section}{0pt}{1.6em}{0.8em}

% --- Header / Footer ---
\usepackage{fancyhdr}
\pagestyle{fancy}
\fancyhf{}
\renewcommand{\headrulewidth}{0.4pt}

% Customise these three commands in your document after % =============================================================
%  preamble.tex — Reusable preamble for Neural Networks course
%  Usage: \input{preamble} at the top of your document,
%         after \documentclass but before \begin{document}.
% =============================================================

% --- Encoding & language ---
\usepackage[T1]{fontenc}
\usepackage[utf8]{inputenc}
\usepackage[english]{babel}         % swap for [french] if needed

% --- Page geometry ---
\usepackage{geometry}
\geometry{top=2.5cm, bottom=2.5cm, left=2.5cm, right=2.5cm}

% --- Math ---
\usepackage{amsmath, amssymb}

% --- Tables ---
\usepackage{booktabs}               % \toprule, \midrule, \bottomrule
\usepackage{array}                  % column type extensions
\usepackage{tabularx}               % auto-width columns
\usepackage{multirow}               % multi-row cells

% --- Colors ---
\usepackage{xcolor}

\definecolor{mainblue}{RGB}{30, 80, 160}
\definecolor{lightblue}{RGB}{220, 235, 255}
\definecolor{accentgray}{RGB}{245, 246, 248}
\definecolor{darkgray}{RGB}{80, 80, 80}
\definecolor{solutiongreen}{RGB}{20, 120, 60}
\definecolor{lightgreen}{RGB}{220, 245, 230}
\definecolor{warnorange}{RGB}{210, 120, 20}
\definecolor{lightorange}{RGB}{255, 243, 220}

% --- Lists ---
\usepackage{enumitem}

% --- TikZ & pgfplots ---
\usepackage{tikz}
\usepackage{pgfplots}
\pgfplotsset{compat=1.18}
\usetikzlibrary{positioning, arrows.meta, calc, patterns}

% --- Boxes: tcolorbox ---
\usepackage{tcolorbox}
\tcbuselibrary{skins, breakable}

% Generic exercise box (blue border, gray background)
\tcbset{
  exercisebox/.style={
    enhanced,
    colback=accentgray,
    colframe=mainblue,
    boxrule=0.8pt,
    arc=4pt,
    left=10pt, right=10pt, top=8pt, bottom=8pt,
    fonttitle=\bfseries\color{mainblue},
  }
}

% Solution box (green, used in answer sheets)
\tcbset{
  solutionbox/.style={
    enhanced,
    colback=lightgreen,
    colframe=solutiongreen,
    boxrule=0.8pt,
    arc=4pt,
    left=10pt, right=10pt, top=8pt, bottom=8pt,
    fonttitle=\bfseries\color{solutiongreen},
  }
}

% Answer highlight box (orange, for final numeric answers)
\tcbset{
  answerbox/.style={
    enhanced,
    colback=lightorange,
    colframe=warnorange,
    boxrule=0.7pt,
    arc=4pt,
    left=8pt, right=8pt, top=5pt, bottom=5pt,
    fonttitle=\bfseries\color{warnorange},
  }
}

% Step box (light blue border, gray background, for intermediate steps)
\tcbset{
  stepbox/.style={
    enhanced,
    colback=accentgray,
    colframe=mainblue!40,
    boxrule=0.5pt,
    arc=3pt,
    left=8pt, right=8pt, top=6pt, bottom=6pt,
  }
}

% --- Inline note / warning box (mdframed) ---
\usepackage{mdframed}
\newmdenv[
  linecolor=warnorange,
  backgroundcolor=orange!5,
  roundcorner=4pt,
  innerleftmargin=8pt,
  innerrightmargin=8pt,
  innertopmargin=6pt,
  innerbottommargin=6pt,
  linewidth=0.7pt
]{notebox}

% --- Section formatting ---
\usepackage{titlesec}

% Section style: "Exercise N : Title" with a blue rule underneath
% Change the prefix via \renewcommand{\sectionprefix}{...} before \section
\newcommand{\sectionprefix}{Exercise}
\titleformat{\section}[block]
  {\large\bfseries\color{mainblue}}
  {\sectionprefix~\thesection\,:}{0.6em}{}
  [\vspace{-0.3em}\textcolor{mainblue}{\rule{\linewidth}{0.8pt}}]
\titlespacing{\section}{0pt}{1.6em}{0.8em}

% --- Header / Footer ---
\usepackage{fancyhdr}
\pagestyle{fancy}
\fancyhf{}
\renewcommand{\headrulewidth}{0.4pt}

% Customise these three commands in your document after \input{preamble}:
%   \renewcommand{\doccoursename}{My Course}
%   \renewcommand{\docsessionname}{Session N — Topic}
\newcommand{\doccoursename}{Neural Networks for Engineers}
\newcommand{\docsessionname}{Session}

\fancyhead[L]{\small\textcolor{darkgray}{\doccoursename}}
\fancyhead[R]{\small\textcolor{darkgray}{\docsessionname}}
\fancyfoot[C]{\small\textcolor{darkgray}{\thepage}}

% --- Title block macro ---
% Usage: \maketitleblock{Title line}{Subtitle / metadata line}{background color}
% Examples:
%   \maketitleblock{Session 4 — Exercises}{Mixed level | Neural Networks}{mainblue}
%   \maketitleblock{Session 4 — Solutions}{Perceptrons and MLPs | Neural Networks}{solutiongreen}
\newcommand{\maketitleblock}[3]{%
  \begin{tcolorbox}[
    enhanced,
    colback=#3,
    colframe=#3,
    arc=6pt, boxrule=0pt,
    left=16pt, right=16pt, top=12pt, bottom=12pt
  ]
    {\large\bfseries\color{white} #1}\\[4pt]
    {\small\color{white!85!black} #2}
  \end{tcolorbox}
  \vspace{1em}
}

% --- Convenience macros ---

% Blank answer field (inline underline for fill-in tables)
\newcommand{\acell}{\rule{1.8cm}{0.4pt}}

% Step-by-step computation display: \step{label}{math}
% Example: \step{Hidden layer}{z^{(1)} = Wx + b = \ldots}
\newcommand{\step}[2]{%
  \noindent\textbf{#1:}\quad $#2$\par\vspace{0.3em}
}

% Neuron style for TikZ diagrams (reuse with \node[neuron])
% Colors are overridable per layer — see examples below.
%
% Input neuron  : fill=lightblue,  draw=mainblue
% Hidden neuron : fill=green!15,   draw=green!50!black
% Output neuron : fill=orange!15,  draw=orange!70!black
%
% Example architecture diagram:
%
%   \begin{tikzpicture}[
%     neuron/.style={circle, minimum size=1cm, thick, font=\small\bfseries},
%     arr/.style={->, >=Stealth, thick, gray!70}
%   ]
%     \node[neuron, fill=lightblue,  draw=mainblue]          (x1) at (0, 1)  {$x_1$};
%     \node[neuron, fill=green!15,   draw=green!50!black]    (h1) at (3, 1)  {$h_1$};
%     \node[neuron, fill=orange!15,  draw=orange!70!black]   (y)  at (6, 1)  {$\hat{y}$};
%     \draw[arr] (x1) -- (h1);
%     \draw[arr] (h1) -- (y);
%   \end{tikzpicture}
:
%   \renewcommand{\doccoursename}{My Course}
%   \renewcommand{\docsessionname}{Session N — Topic}
\newcommand{\doccoursename}{Neural Networks for Engineers}
\newcommand{\docsessionname}{Session}

\fancyhead[L]{\small\textcolor{darkgray}{\doccoursename}}
\fancyhead[R]{\small\textcolor{darkgray}{\docsessionname}}
\fancyfoot[C]{\small\textcolor{darkgray}{\thepage}}

% --- Title block macro ---
% Usage: \maketitleblock{Title line}{Subtitle / metadata line}{background color}
% Examples:
%   \maketitleblock{Session 4 — Exercises}{Mixed level | Neural Networks}{mainblue}
%   \maketitleblock{Session 4 — Solutions}{Perceptrons and MLPs | Neural Networks}{solutiongreen}
\newcommand{\maketitleblock}[3]{%
  \begin{tcolorbox}[
    enhanced,
    colback=#3,
    colframe=#3,
    arc=6pt, boxrule=0pt,
    left=16pt, right=16pt, top=12pt, bottom=12pt
  ]
    {\large\bfseries\color{white} #1}\\[4pt]
    {\small\color{white!85!black} #2}
  \end{tcolorbox}
  \vspace{1em}
}

% --- Convenience macros ---

% Blank answer field (inline underline for fill-in tables)
\newcommand{\acell}{\rule{1.8cm}{0.4pt}}

% Step-by-step computation display: \step{label}{math}
% Example: \step{Hidden layer}{z^{(1)} = Wx + b = \ldots}
\newcommand{\step}[2]{%
  \noindent\textbf{#1:}\quad $#2$\par\vspace{0.3em}
}

% Neuron style for TikZ diagrams (reuse with \node[neuron])
% Colors are overridable per layer — see examples below.
%
% Input neuron  : fill=lightblue,  draw=mainblue
% Hidden neuron : fill=green!15,   draw=green!50!black
% Output neuron : fill=orange!15,  draw=orange!70!black
%
% Example architecture diagram:
%
%   \begin{tikzpicture}[
%     neuron/.style={circle, minimum size=1cm, thick, font=\small\bfseries},
%     arr/.style={->, >=Stealth, thick, gray!70}
%   ]
%     \node[neuron, fill=lightblue,  draw=mainblue]          (x1) at (0, 1)  {$x_1$};
%     \node[neuron, fill=green!15,   draw=green!50!black]    (h1) at (3, 1)  {$h_1$};
%     \node[neuron, fill=orange!15,  draw=orange!70!black]   (y)  at (6, 1)  {$\hat{y}$};
%     \draw[arr] (x1) -- (h1);
%     \draw[arr] (h1) -- (y);
%   \end{tikzpicture}
:
%   \renewcommand{\doccoursename}{My Course}
%   \renewcommand{\docsessionname}{Session N — Topic}
\newcommand{\doccoursename}{Neural Networks for Engineers}
\newcommand{\docsessionname}{Session}

\fancyhead[L]{\small\textcolor{darkgray}{\doccoursename}}
\fancyhead[R]{\small\textcolor{darkgray}{\docsessionname}}
\fancyfoot[C]{\small\textcolor{darkgray}{\thepage}}

% --- Title block macro ---
% Usage: \maketitleblock{Title line}{Subtitle / metadata line}{background color}
% Examples:
%   \maketitleblock{Session 4 — Exercises}{Mixed level | Neural Networks}{mainblue}
%   \maketitleblock{Session 4 — Solutions}{Perceptrons and MLPs | Neural Networks}{solutiongreen}
\newcommand{\maketitleblock}[3]{%
  \begin{tcolorbox}[
    enhanced,
    colback=#3,
    colframe=#3,
    arc=6pt, boxrule=0pt,
    left=16pt, right=16pt, top=12pt, bottom=12pt
  ]
    {\large\bfseries\color{white} #1}\\[4pt]
    {\small\color{white!85!black} #2}
  \end{tcolorbox}
  \vspace{1em}
}

% --- Convenience macros ---

% Blank answer field (inline underline for fill-in tables)
\newcommand{\acell}{\rule{1.8cm}{0.4pt}}

% Step-by-step computation display: \step{label}{math}
% Example: \step{Hidden layer}{z^{(1)} = Wx + b = \ldots}
\newcommand{\step}[2]{%
  \noindent\textbf{#1:}\quad $#2$\par\vspace{0.3em}
}

% Neuron style for TikZ diagrams (reuse with \node[neuron])
% Colors are overridable per layer — see examples below.
%
% Input neuron  : fill=lightblue,  draw=mainblue
% Hidden neuron : fill=green!15,   draw=green!50!black
% Output neuron : fill=orange!15,  draw=orange!70!black
%
% Example architecture diagram:
%
%   \begin{tikzpicture}[
%     neuron/.style={circle, minimum size=1cm, thick, font=\small\bfseries},
%     arr/.style={->, >=Stealth, thick, gray!70}
%   ]
%     \node[neuron, fill=lightblue,  draw=mainblue]          (x1) at (0, 1)  {$x_1$};
%     \node[neuron, fill=green!15,   draw=green!50!black]    (h1) at (3, 1)  {$h_1$};
%     \node[neuron, fill=orange!15,  draw=orange!70!black]   (y)  at (6, 1)  {$\hat{y}$};
%     \draw[arr] (x1) -- (h1);
%     \draw[arr] (h1) -- (y);
%   \end{tikzpicture}
:
%   \renewcommand{\doccoursename}{My Course}
%   \renewcommand{\docsessionname}{Session N — Topic}
\newcommand{\doccoursename}{Neural Networks for Engineers}
\newcommand{\docsessionname}{Session}

\fancyhead[L]{\small\textcolor{darkgray}{\doccoursename}}
\fancyhead[R]{\small\textcolor{darkgray}{\docsessionname}}
\fancyfoot[C]{\small\textcolor{darkgray}{\thepage}}

% --- Title block macro ---
% Usage: \maketitleblock{Title line}{Subtitle / metadata line}{background color}
% Examples:
%   \maketitleblock{Session 4 — Exercises}{Mixed level | Neural Networks}{mainblue}
%   \maketitleblock{Session 4 — Solutions}{Perceptrons and MLPs | Neural Networks}{solutiongreen}
\newcommand{\maketitleblock}[3]{%
  \begin{tcolorbox}[
    enhanced,
    colback=#3,
    colframe=#3,
    arc=6pt, boxrule=0pt,
    left=16pt, right=16pt, top=12pt, bottom=12pt
  ]
    {\large\bfseries\color{white} #1}\\[4pt]
    {\small\color{white!85!black} #2}
  \end{tcolorbox}
  \vspace{1em}
}

% --- Convenience macros ---

% Blank answer field (inline underline for fill-in tables)
\newcommand{\acell}{\rule{1.8cm}{0.4pt}}

% Step-by-step computation display: \step{label}{math}
% Example: \step{Hidden layer}{z^{(1)} = Wx + b = \ldots}
\newcommand{\step}[2]{%
  \noindent\textbf{#1:}\quad $#2$\par\vspace{0.3em}
}

% Neuron style for TikZ diagrams (reuse with \node[neuron])
% Colors are overridable per layer — see examples below.
%
% Input neuron  : fill=lightblue,  draw=mainblue
% Hidden neuron : fill=green!15,   draw=green!50!black
% Output neuron : fill=orange!15,  draw=orange!70!black
%
% Example architecture diagram:
%
%   \begin{tikzpicture}[
%     neuron/.style={circle, minimum size=1cm, thick, font=\small\bfseries},
%     arr/.style={->, >=Stealth, thick, gray!70}
%   ]
%     \node[neuron, fill=lightblue,  draw=mainblue]          (x1) at (0, 1)  {$x_1$};
%     \node[neuron, fill=green!15,   draw=green!50!black]    (h1) at (3, 1)  {$h_1$};
%     \node[neuron, fill=orange!15,  draw=orange!70!black]   (y)  at (6, 1)  {$\hat{y}$};
%     \draw[arr] (x1) -- (h1);
%     \draw[arr] (h1) -- (y);
%   \end{tikzpicture}
